% concatenated from arxiv.tex
\documentclass[11pt, reqno]{amsart}
\usepackage{amsmath,amssymb}
\usepackage{iftex}
\usepackage{tabularray}
%\usepackage{newtxtext,newtxmath}
\usepackage{subcaption}
\usepackage{amsthm,amsmath,amssymb,amsfonts}
\usepackage{natbib}
\usepackage{txfonts}
%\usepackage{fullpage}
\usepackage{booktabs}
\usepackage{float}
\usepackage{graphicx}
\usepackage{graphics}
\usepackage{ascmac}
\usepackage{bm}
\usepackage{accents}
\usepackage[colorlinks=true, allcolors=blue]{hyperref}
%\usepackage{pxjahyper}
\usepackage{url}
\usepackage{algorithm}
\usepackage{algorithmic}
\usepackage{enumitem}
\renewcommand{\thefootnote}{*\arabic{footnote}}


\newcommand{\R}{\mathbb{R}}
\newcommand{\Q}{\mathbb{Q}}
\newcommand{\N}{\mathbb{N}}
\newcommand{\Z}{\mathbb{Z}}
\newcommand{\bF}{\mathbb{F}}
\newcommand{\bG}{\mathbb{G}}
\newcommand{\bD}{\mathbb{D}}
\newcommand{\bE}{\mathbb{E}}
\newcommand{\bP}{\mathbb{P}}
\newcommand{\tr}{^{\top}}

\renewcommand{\tilde}{\widetilde}
\renewcommand{\hat}{\widehat}
\renewcommand{\le}{\leqslant}
\renewcommand{\ge}{\geqslant}

\newcommand{\mA}{\mathcal{A}}
\newcommand{\mB}{\mathcal{B}}
\newcommand{\mC}{\mathcal{C}}
\newcommand{\mD}{\mathcal{D}}
\newcommand{\mF}{\mathcal{F}}
\newcommand{\mG}{\mathcal{G}}
\newcommand{\mH}{\mathcal{H}}
\newcommand{\mK}{\mathcal{K}}
\newcommand{\mL}{\mathcal{L}}
\newcommand{\mN}{\Normal}
\newcommand{\mO}{O}
\newcommand{\mP}{\mathcal{P}}
\newcommand{\mS}{\mathcal{S}}
\newcommand{\mT}{\mathcal{T}}
\newcommand{\mV}{\mathcal{V}}
\newcommand{\mX}{\mathcal{X}}
\newcommand{\mY}{\mathcal{Y}}
\renewcommand{\Re}{\mathrm{Re}}
\newcommand{\indep}{\mathop{\perp\perp}}

\def\vector#1{\mbox{\boldmath $#1$}}

\newcommand{\cvp}{\stackrel{\bP}{\rightarrow}}
\newcommand{\cvd}{\stackrel{d}{\rightarrow}}
\newcommand{\wcv}{\stackrel{w}{\rightarrow}}
\newcommand{\deq}{\stackrel{d}{=}}
\newcommand{\1}{\mbox{1}\hspace{-0.25em}\mbox{l}}

\usepackage{color}
\newcommand{\rc}{\color{red}}
\newcommand{\bc}{\color{blue}}
\newcommand{\gc}{\color{green}}
\usepackage{comment}

\DeclareMathOperator{\Var}{Var}
\DeclareMathOperator{\Cov}{Cov}
\DeclareMathOperator{\Corr}{Corr}
\DeclareMathOperator{\const}{const.}
\newcommand{\Varp}[1]{\Var\left(#1\right)}
\newcommand{\Covp}[1]{\Cov\left(#1\right)}
\newcommand{\Corrp}[1]{\Corr\left(#1\right)}
\newcommand{\WP}{\textrm{ w.p.\ }}
\newcommand{\obs}{\mathrm{obs}}
\newcommand{\rel}{\mathrm{rel}}
\newcommand{\rank}{\mathrm{rank}}
\DeclareMathOperator{\sgn}{sgn}
\DeclareMathOperator{\range}{range}
\DeclareMathOperator*{\argmax}{argmax}
\DeclareMathOperator*{\argmin}{argmin}
\DeclareMathOperator*{\plim}{plim}

\newcommand{\Ind}[1]{\1\left\{#1\right\}}
\newcommand{\Normal}{\mathrm{N}}
\newcommand{\Normalp}[2]{\Normal\left(#1,#2\right)}
\newcommand{\BetaDist}{\textrm{Beta}}
\newcommand{\DirDist}{\textrm{Dir}}
\newcommand{\UnifDist}{\textrm{Unif}}
\newcommand{\weaklyto}{\rightsquigarrow}
\newcommand{\diff}{d} 
\newcommand{\pD}[2]{\frac{\partial #1}{\partial #2}}
\newcommand{\pDD}[2]{\frac{\partial^2 #1}{\partial #2^2}} 
\newcommand{\D}[2]{\frac{\diff #1}{\diff #2}}
\newcommand{\DD}[2]{\frac{\diff^2 #1}{\diff #2^2}} 

\theoremstyle{definition}\newtheorem{problem}{Problem}[section]
\theoremstyle{definition}\newtheorem{definition}[problem]{Definition}
\theoremstyle{remark}\newtheorem{assumption}{Assumption}
\renewcommand*{\theassumption}{\Alph{assumption}}
\theoremstyle{remark}\newtheorem{remark}[problem]{Remark}
\theoremstyle{definition}\newtheorem{example}[problem]{Example}
\theoremstyle{plain}\newtheorem{theorem}[problem]{Theorem}
\theoremstyle{plain}\newtheorem{question}[problem]{Question}
\theoremstyle{plain}\newtheorem{lemma}[problem]{Lemma}
\theoremstyle{plain}\newtheorem{proposition}[problem]{Proposition}
\theoremstyle{plain}\newtheorem{corollary}[problem]{Corollary}
\theoremstyle{plain}\newtheorem{conjecture}[problem]{Conjecture}

\keywords{repro samples method, selective inference, permuted regression, shuffled regression, unmatched data, Hungarian algorithms.}

	\title{Finite-Sample Inference for Sparsely Permuted Linear Regression}
    \author[H. Ota]{Hirofumi Ota}
    \address[H. Ota]{Graduate School of Arts and Sciences, The University of Tokyo.}
    \email{hirofumi-ota@g.ecc.u-tokyo.ac.jp}
    \author[M. Imaizumi]{Masaaki Imaizumi}
    \address[M. Imaizumi]{Graduate School of Arts and Sciences, The University of Tokyo. \\RIKEN Center for Advanced Intelligence Project.}
        \email{imaizumi@g.ecc.u-tokyo.ac.jp}
	\date{\today\  (first version)}

	
	
\begin{document}
\maketitle
\begin{abstract}
We study a linear observation model with an unknown permutation called \textit{permuted/shuffled linear regression}, where responses and covariates are mismatched and the permutation forms a discrete, factorial-size parameter. 
The permutation is a key component of the data-generating process, yet its statistical investigation remains challenging due to its discrete nature.
We develop a general statistical inference framework on the permutation and regression coefficients. First, we introduce a localization step that reduces the permutation space to a small candidate set building on recent advances in the repro samples method, whose miscoverage decays polynomially with the number of Monte Carlo samples. Then, based on this localized set, we provide statistical inference procedures: a conditional Monte Carlo test of permutation structures with valid finite-sample Type-I error control. We also develop coefficient inference that remains valid under alignment uncertainty of permutations. For computational purposes, we develop a linear assignment problem computable in polynomial time and demonstrate that, with high probability, the solution is equivalent to that of the conventional least squares with large computational cost. Extensions to partially permuted designs and ridge regularization are further discussed. Extensive simulations and an application to air-quality data corroborate finite-sample validity, strong power to detect mismatches, and practical scalability.
\end{abstract}




\section{Introduction}

\subsection{Background and motivation}

Modern data integration and privacy mechanisms often disrupt the one-to-one correspondence between covariates and responses.  
In this context, the \emph{shuffled/permuted linear regression model} \citep{abid2017linear, pananjady2017linear} serves as a fundamental regression framework that explicitly accounts for possible mismatches through an unknown permutation acting on the design or response indices.  
This model provides a canonical baseline for studying statistical inference under alignment uncertainty and has become a central object in the recent literature on regression with unobserved correspondences.

Such structures arise in settings involving anonymized or aggregated data \citep{narayanan2008robust}, probabilistic record linkage \citep{fellegi1969theory, christen2012evaluation, winkler2014matching}, and privacy-preserving or decentralized data collection protocols \citep{vatsalan2017privacy, dwork2014algorithmic}.  
In these situations, the pairings between predictor and response units may be only partially known or implicitly inferred, thereby motivating the study of regression models under alignment uncertainty.
We focus on the regime in which only a small fraction of the sample pairs are mismatched, often referred to as \emph{sparsely permuted} data \citep{slawski2019linear}.  
Even a small number of mismatches can distort parameter estimates and invalidate uncertainty quantification if ignored.

The appearance of sparsely permuted data motivates two core inferential questions: \textit{whether any permutation is present in the data, and how to construct confidence sets for the regression coefficients that remain valid under alignment uncertainty}.
Addressing these problems presents methodological challenges, because the permutation is discrete and combinatorial, and the search space grows factorially with the sample size, which undermines local asymptotics, smooth likelihood expansions, and tractable selective conditioning. We therefore adopt a nonasymptotic, finite-sample approach based on the \emph{repro samples method} \citep{xie2022repro, xie2024repro, wang2022finite} and the principle of \emph{locally simultaneous inference} \citep{zrnic2024locally}. At a high level, we localize the permutation space to a data-dependent candidate set and perform inference uniformly over that set; details are presented in Section \ref{sec:methodology}-\ref{sec:inference}.


\subsection{Problem statement}
We consider the linear regression problem with permuted data, where the observed response vector $\bm{Y}^{\obs} \in \R^n$ is paired with an $n \times p$ deterministic design matrix $\bm{X}$ with full column rank $\rank(\bm{X})=p$. The model can be described as 
\begin{equation}\label{eq:model intro}
    \bm{Y}^{\obs}  = \Pi_0 \bm{X} \bm{\beta}_{0} + \sigma_0 \bm{u}^{\rel},
\end{equation}
where $\bm{\beta}_{0} \in \R^p$ is an unknown coefficient vector, $\Pi_0$ is an unknown $n \times n$ permutation matrix, and $\bm{u}^{\rel}$ is a realization from $\bm{U} \sim \Normal(0, I_n)$, which is the noise vector associated with the observation $\bm{Y}^{\obs}$. The scale parameter $\sigma_0 > 0$ is also unknown. Under the frequentist framework, the statistical model in \eqref{eq:model intro} can also be described as the following random-sample version of the data-generating process
\begin{equation}\label{eq:random}
    \bm{Y}  = \Pi_0 \bm{X} \bm{\beta}_{0} + \sigma_0 \bm{U}, \qquad \bm{U} \sim \Normal(0, I_n).
\end{equation}


We further assume that the number of mismatches is at most $k$. Such situation is referred to as sparsely permuted data first appeared in \cite{slawski2019linear}, indicating that the permutation affects only a small fraction of all response-predictor pairs. 
To quantify the number of mismatches, 
we define the Hamming distance between an identity matrix $I_n$ and a $n \times n$ permutation matrix as 
\begin{equation}
d(\Pi, I_n) = \sum_{i = 1}^{n} \1\{\Pi_{ii} = 0\},
\end{equation}
where $\Pi_{ii}$ denotes the $i$-th diagonal element of $\Pi$, and focus on the $k$-sparse permutation matrix that satisfies $d(\Pi, I_n) \le k$. Furthermore, we denote the set of $k$-sparse permutation matrices as 
\begin{equation}
    \mP_{n,k}=\{\Pi\in\mP_n:d(\Pi,I_n)\le k\}.
\end{equation}


Based on the canonical setting, this paper addresses two goals of the statistical inference:
\begin{itemize}
\item[(G1)] \textbf{Testing for mismatches.}  
Construct a finite-sample valid test of the composite hypothesis $H_0: d(\Pi_0,I_n)\le k_0$ (equivalently, $\Pi_0\in\mP_{n,k_0}$) versus $H_1: d(\Pi_0,I_n)>k_0$ (equivalently, $\Pi_0\in\mP_{n,k}\setminus\mP_{n,k_0}$). The special case $k_0=0$ tests whether any mismatch is present.
\item[(G2)] \textbf{Confidence sets for $\bm{\beta}_0$ under unknown $\Pi_0$.}  
Build a confidence set for $\bm{\beta}_0$ that accounts for the uncertainty in estimating $\Pi_0$ and attains nominal coverage in finite samples.
\end{itemize}




\subsection{Overview of contributions}
We develop a methodology for statistical inference toward the above goals, which consists of two stages: generating candidate sets and conducting inference based on them.
In the first step, we generate a candidate set of plausible permutations based on the repro samples method \citep{xie2022repro,xie2024repro}. We generate artificial noise $L$ times, and for each noise draw, we optimize an objective function that mimics the least-squares objective under the corresponding noise realization, and then collect the resulting optimizers into a data-dependent candidate set $\mathcal{C}^{(L)}\subset\mathcal{P}_{n,k}$. This localized domain is then used to conduct coefficient inference by taking union of projection-based $F$-statistics over $\mathcal{C}^{(L)}$, which yields locally simultaneous confidence regions for $\bm{\beta}_0$.
In the second step, we develop a test for the hypothesis about the permutation structure: $H_0:\Pi_0\in\mP_{n,k_0}\quad\text{vs.}\quad H_1:\Pi_0\in\mP_{n,k}\setminus\mP_{n,k_0}$. We construct a discrete statistic based on $\mC^{(L)}$ and calibrate it via conditional Monte Carlo uniformly over the localized null set $\hat{\mP}_{n,k_0}:=\mC^{(L)}\cap\mP_{n,k_0}$. These steps achieve the two statistical inference goals (G1) and (G2).

We summarize the contributions of our study.
\begin{itemize}
    \item Our developed approach is theoretically guaranteed to be valid. For the generation of the candidate set in the first step with the localization, we prove a polynomial-rate bound on the miscoverage probability $\mathbb{P}(\Pi_0\notin\mathcal{C}^{(L)})$ in terms of the number of repro samples $L$ (Theorem~\ref{thm:recovery}). Also, we show the coverage of the confidence regions for $\bm{\beta}_0$ at the nominal level up to the localization error that vanishes at a polynomial rate in $L$ (Theorem~\ref{thm:confidence}). For testing, the developed statistical test controls type-I error at any sample size (Theorem~\ref{thm:p-value}) localized on the event $\{\Pi_0\in\mC^{(L)}\}$. Marginally, the size inflation is bounded by the miscoverage probability of $\mC^{(L)}$.
    \item We develop a computationally effective scheme to implement the method. In particular, we develop a score-weighted linear assignment problem (LAP) as a surrogate to the original least square problem for generating the candidate set. This surrogate problem is solved by the Hungarian algorithm with the computational complexity $O(n^3)$, which avoids the combinatorial explosion of the computational complexity of the original problem. We also prove the high-probability equivalence of the surrogate problem to the original problem (Theorem~\ref{thm:high-probability equivalence}). 
    \item Our methodology is comprehensively validated through both simulation and real-data analysis in an applicati   on to Beijing air-quality data. We also apply our approach  to the partially permuted designs and the permuted ridge regression model.
    Further, we propose selection methods for several necessary hyperparameters, and their proper functioning is verified experimentally.
\end{itemize}


\subsection{Related works}

\subsubsection{Shuffled regression}
Several recent methods use Hungarian/LAP solvers as computational surrogates for shuffled regression, e.g., spectral matching via covariance eigenbasis alignment \citep{liu2025shuffled} or LAP subproblems within alternating schemes \citep{li2021generalized}. 
Relatedly, local search over $k$-sparse permutations performs greedy swaps with monotone loss descent but offers no guarantee of reaching a global optimum in noisy settings \citep{mazumder2023linear}. 
In contrast, Theorem~\ref{thm:high-probability equivalence} provides a finite-sample certification: under explicit conditions, our score-weighted LAP attains the same minimizer as the repro-sample objective, and thus agrees with its global minimizer with high probability.


Research related to shuffled regression includes the estimation of shape-constrained models when the correspondence between covariates and responses is unknown.
For one-dimensional covariate, \cite{rigollet2019uncoupled} obtain minimax rates of an estimation error for uncoupled isotonic regression via minimum Wasserstein deconvolution. 
In multivariate settings, there are more studies.
For coordinate-wise isotonic regression with axis-wise permutations, \cite{pananjady2022isotonic} derive minimax and adaptation results and propose a polynomial-time partition estimator. \cite{slawski2024permuted} establish identification via cyclical monotonicity and develop optimal-transport/likelihood-based estimators with certain error bounds. Permutation uncertainty also appears in structured arrays: \cite{ma2021optimal} give information-theoretic limits and near-optimal polynomial-time recovery for monotone matrices, and \cite{lee2025statistical} provide statistically optimal and computationally efficient estimators for smooth tensors. 

A complementary line of work develops statistical inference for shuffled linear regression. \cite{slawski2021pseudo} propose a pseudo-likelihood-based asymptotic inference for the coefficient vector and provide EM-type algorithm for the optimization. \cite{chakraborty2024learning} study sparsely permuted data from a robust Bayesian perspective, treating mispaired entries as outliers within a generalized fractional posterior, and establish posterior consistency alongside practical sampling strategies.


\subsubsection{Repro samples method}


Our work adopts the repro samples method \citep{xie2022repro,xie2024repro}, which provides non-asymptotic, simulation-based inference by quantifying uncertainty through artificial samples. 
In high-dimensional linear models, repro samples have been used to construct model confidence sets and valid inference for coefficients \citep{wang2022finite}, with extensions to high-dimensional logistic models \citep{hou2025repro}. 
We transfer these ideas from support uncertainty (scaling primarily with $p$) to the qualitatively different regime of permutation uncertainty (scaling factorially with $n$): we build a data-dependent candidate set of permutations $\mC^{(L)}$ that contains the truth with high probability, then conduct inference uniformly over $\mC^{(L)}$ in the spirit of locally simultaneous inference. 
This yields (i) a finite-sample valid test for the presence of any mismatch and (ii) union-type confidence sets for $\bm{\beta}_0$ that account for alignment uncertainty. 

\subsection{Organization}
Section~\ref{sec:methodology} develops a localization step via repro samples, yielding a data-dependent candidate set of permutations and a practical score-weighted LAP. Section~\ref{sec:inference} then introduces finite-sample procedures: a conditional Monte Carlo test of no permutation and coefficient inference that remains valid under alignment uncertainty. Section~\ref{sec:theory} presents the theoretical guarantees: oracle identifiability, polynomial-rate miscoverage bounds for the candidate set, finite-sample validity of the test, coverage of the coefficient regions, and a high-probability argmin equivalence for the Hungarian surrogate. Section~\ref{sec:extensions} discusses partially permuted designs and ridge regularization. Section \ref{sec:tuning} discusses how to select the tuning parameters and Section~\ref{sec:numerical} reports empirical evidence from extensive simulations and an application to Beijing air-quality data. Finally, Section~\ref{sec:conclusion} concludes. The Appendix provides methodological background, additional theory, full proofs, and supplementary figures.



\subsection{Notation}
Vectors and matrices are bold.  
$\|\cdot\|_2$ denotes the Euclidean norm, $\|\cdot\|_\infty$ the max norm, and $\|\cdot\|_{\mathrm{op}}$ the operator norm.  
For a matrix $\bm A\in\R^{n\times p}$, $M_{\bm{A}} \in \R^{n \times n}$ denotes the orthogonal projection matrix onto $\range(\bm{A})$, and $R_{\bm{A}}:=I_n-M_{\bm{A}}\in \R^{n \times n}$ denotes the residual projection matrix.
The Hamming distance between permutations $\Pi$ and $I_n$ is $d(\Pi,I_n)=\sum_{i=1}^n\1\{\Pi_{ii}=0\}$.
We use $\mP_n$ for the set of all $n\times n$ permutation matrices and $\mP_{n,k}:=\{\Pi\in\mP_n: d(\Pi,I_n)\le k\}$ for the $k$-sparse subset.  
For random variables, $\Normal(\mu,\Sigma)$ denotes the Gaussian distribution, and $F_{a,b}$ denotes an $F$-distribution with $(a,b)$ degrees of freedom.  


\section{Candidate set of permutations via repro samples}\label{sec:methodology}
We develop the candidate set of permutation matrices, which is the key ingredient of our statistical inference framework.
This section includes (i) a localization approach based on the repro sample method, and (ii) an score-weighted linear assignment problem (LAP) to handle our setup.





\subsection{Localization principle}\label{subsec:candidate}


We present the localization approach to make a candidate set $\mC^{(L)}\subset\mP_{n,k}$, which focuses on a small subset containing the permutation $\Pi_0$ with high probability and avoids the exhaustive search over all permutations in $\mP_{n,k}$.
This approach leverages the fact that the permutation estimated under artificially generated noise, known as a repro sample, is close to the true permutation identified under the realized noise.

To illustrate the idea, we first consider an ideal situation in which the unobserved realization of noise $\bm{u}^{\rel}$ in \eqref{eq:model intro} is fully available. The following Proposition \ref{prop:identification} claims that we can identify the true permutation $\Pi_0$ with $\bm{u}^{\rel}$ under $k$-sparse setup. 


We make the following baseline assumption on the design matrix $\bm X$.
\begin{assumption}[Design matrix]\label{ass:generic-design}
The rows of $\bm X$ are generated i.i.d.\ from a distribution on
$\R^p$ with a Lebesgue density.
\end{assumption}
Throughout the paper, we treat $\bm X$ as a fixed design matrix and
conduct inference conditionally on its realized value.
The assumption above is imposed solely to guarantee that, with
probability one, the realized design matrix satisfies the required
full-rank condition.




\begin{proposition}[Unique recovery in the oracle setting]\label{prop:identification}
    In addition to Assumption \ref{ass:generic-design}, suppose $n-2k \ge p$ and $\bm\beta_0\neq \bm 0$ hold. Then, the true permutation $\Pi_0$ uniquely satisfies 

\begin{equation}\label{eq:repro_ideal_Pi}
   \Pi_0 = \argmin_{\Pi \in \mP_{n,k}}\| (I_n - M_{(\Pi \bm{X}, \bm{u}^{\rel})}) \bm{Y}^{\obs}\|_2^2,
\end{equation}
where $M_{(\Pi \bm{X}, \bm{u}^{\rel})}$ is the orthogonal projector onto the range of $(\Pi \bm{X}, \bm{u}^{\rel})$. 
\end{proposition}

When $\bm{u}^{\rel}$ is available, Proposition \ref{prop:identification} states that $\Pi_0$ is uniquely recovered. Also, we note that the $k$-sparse setup enables this recovery result. The complete proof is deferred to Appendix \ref{proof:prop identification}.


Next, we construct a candidate set in the case where $\bm{u}^{\rel}$ is not available.
By utilizing the insight from the oracle result as Proposition \ref{prop:identification}, we generate artificial samples of the noise that mimic the objective in \eqref{eq:repro_ideal_Pi}. Specifically, we generate $L$ i.i.d.\ vectors $\bm{u}_1^*, \ldots, \bm{u}_L^* \sim \Normal(0, I_n)$ 
and, for each $\ell=1,\ldots,L$, we define the least square problem and its optimizer
\begin{equation}\label{eq:repro_Pi}
   \hat{\Pi}_{\ell} = \argmin_{\Pi \in \mP_{n,k}}\| (I_n - M_{(\Pi \bm{X}, \bm{u}_{\ell}^*)}) \bm{Y}^{\obs}\|_2^2.
\end{equation}
This is an analogue of the oracle problem \eqref{eq:repro_ideal_Pi} by replacing the realization $\bm{u}^{\rel}$ by its artificial counterpart $\bm{u}^*_\ell$. We expect $\hat{\Pi}_{\ell}$ is close to $\Pi_0$ when $\bm{u}_{\ell}^*$ is sufficiently aligned with $\bm{u}^{\rel}$.

Finally, we collect $\hat{\Pi}_{1},\ldots,\hat{\Pi}_{L}$ to form the candidate set
\begin{equation}\label{eq:permutation candidates set}
    \mathcal{C}^{(L)}=\left\{\hat{\Pi}_{\ell}: \ell =1, \ldots, L\right\}.
\end{equation}


\subsection{Score-weighted linear assignment problem} \label{sec:algorithm}


We develop a methodology to approximately solve the least square objective in \eqref{eq:repro_Pi}.
Since the objective in \eqref{eq:repro_Pi} minimizes a nonlinear projection residual over permutation matrices,  the problem remains computationally demanding due to its combinatorial nature even when we restrict the search to $\mP_{n,k}$ (see \cite{pananjady2016linear,pananjady2017linear}).
To obtain a tractable surrogate, we replace \eqref{eq:repro_Pi} by a score-weighted linear assignment problem (LAP), which has the form $\min_{\Pi} \sum_{i} \sum_j \Omega_{ij} \Pi_{ij}$ with a cost matrix $\Omega \in \R^{n \times n}$ and can be solved exactly in $O(n^3)$ time by the Hungarian algorithm \citep{kuhn1955hungarian}. 


To develop the LAP problem as a surrogate, for a fixed repro sample $\bm u^*$, we reform the original objective \eqref{eq:repro_Pi} into a (scaled) LAP cost plus a projection-mismatch term. Concretely, we expand \eqref{eq:repro_Pi} as follows: \begin{align} 
&\left\|(I_n-M_{(\Pi\bm X,\bm u^*)})\bm Y^{\obs}\right\|_2^2 \label{eq:lap-core-decomp}\\
&=
\underbrace{\frac{1}{2}\sum_{i=1}^n\sum_{j=1}^n (Y_i-m_j)^2\Pi_{ij}}_{=:T_1(\Pi)}
+ \underbrace{\left(\sum_{i,j} Y_i m_j\Pi_{ij}-{\bm Y^{\obs}}^\top M_{(\Pi\bm X,\bm u^*)}\bm Y^{\obs}\right)}_{=: T_2(\Pi)}
+ \underbrace{\frac12\left(\|\bm Y^{\obs}\|_2^2-\sum_{j=1}^n m_j^2\right)}_{=:T_3}. \notag 
\end{align}
The first term $T_1(\Pi)$ is a standard LAP objective with costs $(Y_i-m_j)^2$ where $m_j$ denotes the $j$-th element of $M_{(\bm{X}, \bm{u}^*)}\bm{Y}^{\obs}$, and the third term $T_3$ is ignorable since it does not depend on $\Pi$.
About the second term $T_2(\Pi)$, since $I_n$ is one of the minimizers of the term, i.e., $T_2(I_n) = \min_{\Pi \in \mP_n} T_2(\Pi)$, we can approximate the minimization of $T_2(\Pi)$ by augmenting a cost which induces an identity to be a minimizer. 
In particular, with some coefficients $\lambda_1, \lambda_2 > 0$, we consider a LAP with a sparsity-inducing cost $\lambda_1\1\{j\neq i\} -\lambda_2\1\{j=i\}(Y_i-m_i)^2$, where the first term penalizes each off-diagonal reassignment (an $\ell_0$-type cost on the number of moved indices), and the second term adds a credit to the diagonal residual. % {\rc [We have relaxed the constraints, but note that sparseness with respect to k is incorporated into the LAP cost function.]}


Consequently, we consider an estimator by solving the following score-weighted LAP as a surrogate for the least square problem \eqref{eq:repro_Pi}
\begin{align}  \label{def:tilde_Pi}
  \tilde{\Pi}
  = \argmin_{\Pi \in \mP_n}
    \sum_{i=1}^n \sum_{j=1}^n \Omega_{ij}(\bm{u}^*)\Pi_{ij},
\end{align}
where $\Omega(\bm{u}^*) \in \R^{n \times n}$ is a cost matrix with $\bm{u}^*$ defined as
\begin{align} \label{eq:Omega_ustar}
  \Omega_{ij}(\bm{u}^*)
  =  (Y_i - m_j)^2 + \lambda_1\1\{j\neq i\}
     -\lambda_2 \1\{j=i\}(Y_i - m_i)^2.
\end{align}
Theorem~\ref{thm:high-probability equivalence} formalizes this reduction: for suitable $(\lambda_1,\lambda_2)$, the score-weighted LAP minimizer coincides with the argmin of \eqref{eq:repro_Pi} with high probability.
We will further discuss the selection of the parameters $\lambda_1$ and $\lambda_2$ in Section \ref{sec:tuning}.
For the sake of completeness, we present the original Hungarian algorithm in Algorithm \ref{alg:scoring hungarian} in Appendix \ref{sec:hungarian alg}.

\subsection{Unique recovery of the parameters}

We extend the results of the recovery of permutation matrices $\Pi_0$ in Proposition \ref{prop:identification} to include coefficients $(\bm{\beta}_0, \sigma_0) $. This is crucial for performing parameter inference.

\begin{proposition}\label{prop:identification_joint}
    If $n-2k\ge p$ holds, then, given the noise realization $\bm u^{\rel}$, the triple $(\Pi_0,\bm\beta_0,\sigma_0) \in \mP_{n,k} \times \R^p \times \R_+$ is the unique minimizer of
\begin{equation}\label{eq:repro_ideal}
   \min_{\Pi \in \mP_{n,k},\ \bm{\beta} \in \R^p,\ \sigma \ge 0 }
   \left\| \bm{Y}^{\obs} -  \Pi \bm{X} \bm{\beta} - \sigma \bm{u}^{\rel}\right\|_2^2.
\end{equation}
\end{proposition}
 
We compare the identifiability condition $n-2k\ge p$ in Proposition \ref{prop:identification_joint} to that of a previous study without the $k$-sparsity.
In particular, for the recovery problem without any constraints,  \citep{unnikrishnan2018unlabeled} derived the identifiability condition $n\ge 2p$.
Compared with this condition, our derived condition shows that $k$-sparsity can reduce the required sample size compared to the unrestricted case when $k$ is small.
See Theorem \ref{thm:counterexample} in Appendix \ref{section:appendix theory} for details.


\section{Finite-sample statistical inference}\label{sec:inference}

\subsection{Inference on the permutation matrix}\label{subsubsection:testing permutation}
This section construct a statistic for the $k_0$-sparsity test with the hypotheses
\begin{equation}\label{eq:hypothesis-k0}
    H_0:\Pi_0\in\mP_{n,k_0}\quad\text{vs.}\quad H_1:\Pi_0\in\mP_{n,k}\setminus\mP_{n,k_0},
\end{equation}
where $k_0 \le k$, by introducing a conditional Monte Carlo method for valid finite-sample statistical inference of the permutation matrix $\Pi_0$. The special case with $k_0=0$ reduces to testing whether any mismatch is present (i.e., $\Pi_0 = I_n$). In the context of shuffled regression, the presence of a permutation is of direct interest \citep{slawski2021pseudo}. 

\subsubsection{Applying localization}


First, we localize the null domain using the candidate set $\mC^{(L)}$ constructed in Section \ref{sec:methodology}.
The key property established there is that $\mC^{(L)}$ contains the true permutation $\Pi_0$ with high probability (see Theorem \ref{thm:recovery}).
Localization restricts attention to the data-dependent but typically small set $\hat{\mP}_{n,k_0}:=\mC^{(L)}\cap\mP_{n,k_0}$, on which we calibrate the rejection threshold uniformly over $\Pi\in\hat{\mP}_{n,k_0}$.
This approach reduces the computational burden of the conditional Monte Carlo calibration under
$H_0:\Pi_0\in\mP_{n,k_0}$, which requires controlling null distribution whose size is combinatorial in the sample size $n$.

More formally, we compute a least squares estimator over the candidate set,
\begin{equation}\label{eq:Pi-hat-LS}
    \hat{\Pi}(\bm Y)\in\argmin_{\Pi\in\mC^{(L)}}\left\|(I_n-M_{\Pi\bm X})\bm Y\right\|_2^2,
\end{equation}
and use the test statistic
\begin{equation}\label{eq:Dk-stat}
    D(\bm Y):=d\left(\hat{\Pi}(\bm Y),I_n\right),\qquad D^{\obs}:=D(\bm Y^{\obs}).
\end{equation}
Intuitively, under $H_0$, the best-fitting permutation in $\mC^{(L)}$ should typically have at most $k_0$ mismatches, whereas under $H_1$, it must typically use more than $k_0$ mismatches to explain the data. 

\subsubsection{Handling nuisance}

Next, we remove an effect of the nuisance parameter $(\bm{\beta}, \sigma)$ to enable inference focused on $\Pi$, which is our primary interest.
To this aim, we consider a statistical model $\bm{Y}_{\bm{\theta}} = \Pi \bm{X} \bm{\beta} + \sigma \bm{U}$ with $\bm{\theta} = (\Pi, \bm{\beta}, \sigma) \in \mP_{n,k} \times \R^p \times \R_+$ and $\bm U\sim\Normal(0,I_n)$, then study a conditional distribution of $\bm{Y}_{\bm{\theta}}$, which is necessary for the conditional Monte Carlo approximation.


We derive the conditional distribution of $ \bm{Y}_{\bm \theta}$ using a formula of sufficient statistics.
By the reformulation $
\bm{Y}_{\bm{\theta}} = \bm{M}_{\Pi \bm{X}} \bm{Y}_{\bm{\theta}} + (I_n - \bm{M}_{\Pi \bm{X}})\bm{Y}_{\bm{\theta}} = \bm{M}_{\Pi \bm{X}} \bm{Y}_{\bm{\theta}} + \sigma(I_n - \bm{M}_{\Pi \bm{X}})\bm{U}
$ with a fixed $\Pi$, 
we define a pair of sufficient statistics for $(\bm \beta, \sigma)$ as
\begin{equation}\label{eq:sufficient statistics}
    S_{1,\Pi}(\bm Y_{\bm \theta}):=M_{\Pi\bm X}\bm Y_{\bm \theta},\qquad S_{2,\Pi}(\bm Y_{\bm \theta}):=\left\|(I_n-M_{\Pi\bm X})\bm Y_{\bm \theta}\right\|_2.
\end{equation}
Then, 
we have the representation of the conditional distribution as
\begin{equation}\label{eq:conditional-representation}
    \bm Y_{\bm \theta}  \mid \left(S_{1,\Pi}(\bm Y_{\bm \theta}),S_{2,\Pi}(\bm Y_{\bm \theta})\right)
     \overset{d}{=}\
    S_{1,\Pi}(\bm Y_{\bm \theta})+S_{2,\Pi}(\bm Y_{\bm \theta})\frac{(I_n-M_{\Pi\bm X})\bm U}{\|(I_n-M_{\Pi\bm X})\bm U\|_2}.
\end{equation}
For a fixed $\Pi$, the pair $(S_{1,\Pi}(\bm Y_{\bm \theta}), S_{2,\Pi}(\bm Y_{\bm \theta}))$ is sufficient for $(\bm \beta, \sigma)$, and the direction $\frac{(I_n-M_{\Pi\bm X})\bm U}{\|(I_n-M_{\Pi\bm X})\bm U\|_2}$ is ancillary, being independent of $(\bm \beta, \sigma)$. Consequently, the conditional distribution of $\bm Y_{\bm \theta}  \mid (S_{1,\Pi}(\bm Y_{\bm \theta}),S_{2,\Pi}(\bm Y_{\bm \theta}))$ is free of $(\bm \beta, \sigma)$. This yields the conditional Monte Carlo approximation of the distribution of $D(\bm Y_{\bm \theta})$ in \eqref{eq:Dk-stat}.

\subsubsection{Critical value}

We derive a critical value localized on $\hat{\mP}_{n,k_0}$ by the conditional Monte Carlo approximation.
Given $\bm Y^{\obs}$, we set $s_{1,\Pi}^{\obs}=S_{1,\Pi}(\bm Y^{\obs})$ and $s_{2,\Pi}^{\obs}=S_{2,\Pi}(\bm Y^{\obs})$. 
Let $\widehat c_{1-\alpha}(\Pi)$ denote $(1-\alpha)$-quantile of the distribution of $D(\bm Y) \mid (S_{1,\Pi}(\bm Y),S_{2,\Pi}(\bm Y)) =(s_{1,\Pi}^{\obs}, s_{2,\Pi}^{\obs})$, and define the composite critical value localized on $\hat{\mP}_{n,k_0}$:
\begin{equation}\label{eq:critical-value-composite}
    \widehat c_{1-\alpha}^{(k_0)}:=\max_{\Pi\in\hat{\mP}_{n,k_0}}\widehat c_{1-\alpha}(\Pi), \mbox{~~and~~}
     \hat B^{(k_0)}_{1-\alpha}:=\{0,1,\ldots,\widehat c_{1-\alpha}^{(k_0)}\}.
\end{equation}
We reject $H_0$ when $D^{\obs}\notin \hat B^{(k_0)}_{1-\alpha}$ holds, i.e., when $D^{\obs}>\hat c_{1-\alpha}^{(k_0)}$ holds. To approximate the critical value, we generate i.i.d.\ $\bm u_1^*,\ldots,\bm u_M^*\sim\Normal(0,I_n)$ and define
\[
    \bm Y_{m,\Pi}^*:=s_{1,\Pi}^{\obs}+s_{2,\Pi}^{\obs}\frac{(I_n-M_{\Pi\bm X})\bm u_m^*}{\|(I_n-M_{\Pi\bm X})\bm u_m^*\|_2},
    \qquad m=1,\ldots,M.
\]
For each $m=1,\dots, M$, we compute a test statistic $D_{m,\Pi}^*:=D(\bm Y_{m,\Pi}^*)$ as \eqref{eq:Dk-stat} and use the $(1-\alpha)$-empirical quantile of the conditional Monte Carlo samples $\{ D_{m,\Pi}^*\}_{m=1}^M$.

Theorem \ref{thm:p-value} in Section \ref{sec:theory} shows that, the resulting test controls the type-I error at level $\alpha$ in finite samples conditionally on $s_{1,\Pi}^{\obs}$ and $s_{2,\Pi}^{\obs}$. We also obtain the marginal type-I error control, where the only additional error comes from the miscoverage probability of the candidate set.

\begin{remark}
An alternative test statistic is possible. For example, the statistic $| \min_{\Pi \in \mP_{n,k}}\| (I_n - M_{\Pi \bm{X}}) \bm{Y}_{\bm{\theta}}\|_2^2 - \| (I_n - M_{\bm{X}}) \bm{Y}_{\bm{\theta}}\|_2^2 |$
is continuous and supported on $[0,\infty)$, which avoids computing the minimizer explicitly but complicates conditional Monte Carlo approximation.
\end{remark}



\subsection{Inference on the coefficient vector}\label{subsec:coef}
We now construct a finite-sample valid confidence set for the coefficient vector $\bm{\beta}_0$, accounting for uncertainty from estimating the permutation. Since the parameter of interest is $\bm{\beta}$, we treat $\Pi$ and $\sigma$ as nuisance. We first consider a pivotal map that depends only on $(\Pi, \bm{\beta})$, and then quantify uncertainty about $\Pi$ via the candidate set $\mathcal{C}^{(L)}$. The resulting region is the union of coefficient confidence sets obtained by fixing each $\Pi \in \mathcal{C}^{(L)}$.

We define a statistics for the purpose. With given $(\Pi, \bm{\beta})$, we define it as
\begin{align*}
    T(\bm Y_{\bm \theta}, (\Pi, \bm{\beta}))  
    &= \frac{\bm{U}^\top \bm{M}_{\Pi \bm{X}}\bm{U}/p}{\bm{U}^\top (I_n - \bm{M}_{\Pi \bm{X}}) \bm{U}/(n-p)} 
     = \frac{(\bm Y_{\bm \theta} - \Pi \bm{X} \bm{\beta})^\top \bm{M}_{\Pi \bm{X}} (\bm Y_{\bm \theta} - \Pi \bm{X} \bm{\beta} )/p}{(\bm Y_{\bm \theta} - \Pi \bm{X} \bm{\beta} )^\top (I_n - \bm{M}_{\Pi \bm{X}})(\bm Y_{\bm \theta} - \Pi \bm{X} \bm{\beta} )/(n-p)}.
\end{align*}
Under the case $(\Pi, \bm{\beta}) = (\Pi_0, \bm{\beta}_0)$ with the true parameters, this statistic has an $F_{p, n-p}$ distribution.

Then, we define the level-$\alpha$ confidence set as
\begin{equation}\label{eq:confidence set beta zero}
    \Gamma_{\alpha}^{\mathrm{coef}}(\bm{Y}^{\obs}) = \bigcup_{\Pi \in \mathcal{C}^{(L)}} \left\{\bm{\beta}\in \R^p : T(\bm{Y}^{\obs}, (\Pi, \bm{\beta})) \le F_{p, n-p}^{-1}(\alpha) \right\},
\end{equation}
where $F_{p, n-p}^{-1}$ is the quantile function of the $F_{p, n-p}$ distribution. We show in Theorem \ref{thm:confidence} how the coverage probability of $\Gamma_{\alpha}^{\mathrm{coef}}(\bm{Y}^{\obs})$ is controlled by the nominal level $\alpha$ and by the miscoverage probability $\delta$ of the candidate set.

\section{Theory}\label{sec:theory}


\subsection{High-probability recovery of $k$-sparse permutation matrices}

We study the candidate set \eqref{eq:permutation candidates set} and show that it contains the true permutation $\Pi_0$ with high-probability.
Let $\bP_{\left(\mathcal{U}^L, \bm{U}\right)}$ denote probability with respect to the joint law of $(\mathcal{U}^L, \bm{U})$, where $\mathcal{U}^L = (\bm{U}_1,\dots, \bm{U}_L)$. 

To state the main results, we define the numerical constant
\begin{equation}\label{eq:c min}
     C_{\min}(\Pi_0) = \min_{\substack{\Pi \in \mP_{n,k}\\ \Pi \neq \Pi_0}} \frac{\|(I_n - \Pi M_{\bm{X}} \Pi^\top) \Pi_0\bm{X}\bm{\beta}_{0} \|_2^2}{n \max (d(\Pi_0, I_n) - d(\Pi, I_n), 1)},
\end{equation}
where $M_{\bm{X}} = \bm{X} (\bm{X}^\top \bm{X})^{-1}\bm{X}^\top$. %Our baseline assumption for the main theorem is $C_{\min}(\Pi_0) > 0$.
The constant $C_{\min}(\Pi_0)$ is the degree of separation between the true structure $\Pi_0\bm{X}\bm{\beta}_{0}$ and its projections onto shuffled designs.


We give a theoretical analysis on the miscoverage probability of the candidate set $\mathcal{C}^{(L)}$.




\begin{theorem}[Coverage probability of the candidate set]\label{thm:recovery}
Suppose $C_{\min}(\Pi_0) > 0$ and $n > p + 1$ hold. Then, there exists a positive real sequence $(\delta_L)_{L \in \N}$ satisfying
\begin{equation}
\label{eq:recovery-finite-L}
\bP_{\left(\mathcal{U}^L, \bm{U}\right)}\left( \Pi_0 \notin \mathcal{C}^{(L)} \right) \le
\delta_L,
\end{equation}
for all $L \ge 2$, such that we have the following as $L \to \infty$:
\begin{equation}\label{eq:delta_L}
     \delta_L
= O\left( \frac{\left(\log L\right)^{(n-p)/2}}{L^{(n-p-1)/(2(n-1))}} \right).   
\end{equation}

\end{theorem}






Theorem~\ref{thm:recovery} theoretically validates the localization principle that underlies our framework.  
The candidate set $\mathcal{C}^{(L)}$ acts as an empirical analogue of a confidence region for $\Pi_0$: as $L$ grows, the probability that $\Pi_0$ is excluded from $\mathcal{C}^{(L)}$ decays polynomially.  
Thus, Theorem~\ref{thm:recovery} shows that $\mathcal{C}^{(L)}$ provides a localized representation of the combinatorial parameter space of $\Pi$, turning an intractable search into a manageable and possibly small random subset that contains the truth with high probability.


Theorem~\ref{thm:recovery} is the permutation analogue of the model candidate set results  \citep{wang2022finite,hou2025repro} by the repro-sample approach: whereas those works treat support uncertainty (scaling with $p$), Theorem~\ref{thm:recovery} addresses permutation uncertainty (scaling with $n$), providing the first finite-sample localization guarantee on the discrete setup. 
%for the at this combinatorial scale and extending the repro-sample methodology to a new discrete regime.


\begin{remark}[Comparison with selective inference]
    Theorem~\ref{thm:recovery} offers a principled alternative to classical selective inference, where exact conditioning on the selection event is analytically and computationally infeasible for permutations due to their nonpolyhedral, combinatorial nature.  
Our candidate-set approach replaces exact conditioning with a localization step: it identifies a small data-dependent subset that contains $\Pi_0$ with high probability and then performs inference uniformly over this subset.  
\end{remark}


\subsection{A finite-sample valid test for the permutation matrix}
We now formalize the validity of the $k_0$-sparsity test introduced in Section~\ref{subsubsection:testing permutation}.
Fix $k_0\le k$ and define the statistic $D(\bm Y)$ in \eqref{eq:Dk-stat} and the localized null set $\hat{\mP}_{n,k_0}$.

Then, we validate the critical value $\hat{c}_{1-\alpha}^{(k_0)}$ and the confidence set $\hat{B}_{1-\alpha}^{(k_0)}$ in \eqref{eq:critical-value-composite}.



\begin{theorem}[Type-I error control for the sparsity test] \label{thm:p-value}
Suppose $C_{\min}(\Pi_0) > 0$ and $n > p + 1$ hold.
Then, for the true parameter $\theta_0=(\Pi_0,\bm\beta_0,\sigma_0)$ with $\Pi_0\in\mP_{n,k_0}$, we have the finite-sample bound
\begin{equation}\label{eq:marginal-validity-k0}
    \bP_{\left(\mathcal{U}^L, \bm{U}\right)}\left(D(\bm Y_{\theta_0})\notin \hat{B}^{(k_0)}_{1-\alpha}\right)\le \alpha+\delta_L,
\end{equation}
where $\delta_L$ is specified in Theorem \ref{thm:recovery}.
\end{theorem}

There are two terms in the right-hand side of \eqref{eq:marginal-validity-k0}.
The second term $\delta_L$ is the error from the localization, which is revealed in Theorem \ref{thm:recovery}.
The first term $\alpha$ is the type-I error without the localization. 
To derive the term, we define the composite critical value with fixed $\Pi$:
\begin{equation}\label{eq:c-composite-theory}
    c_{1-\alpha}^{(k_0)}:=\max_{\Pi\in {\mP}_{n,k_0}}c_{1-\alpha}(\Pi), \mbox{~~and~~}
    B^{(k_0)}_{1-\alpha}:=\{0,1,\ldots,c_{1-\alpha}^{(k_0)}\},
\end{equation}
where $c_{1-\alpha}(\Pi)$ denotes the $(1-\alpha)$-quantile of the conditional distribution of $D(\bm Y_{\bm \theta})$ given $(S_{1,\Pi}(\bm Y_{\bm \theta}),S_{2,\Pi}(\bm Y_{\bm \theta}))$, then we show $\bP(D(\bm Y_{\theta_0})\notin B^{(k_0)}_{1-\alpha}\mid S_{1,\Pi_0}(\bm Y_{\theta_0}),\ S_{2,\Pi_0}(\bm Y_{\theta_0}))\le \alpha$ (See Proposition \ref{thm:p-value_without_delta}).
The key ingredients are the conditional representation \eqref{eq:conditional-representation} and the fact that $c_{1-\alpha}^{(k_0)}$ is chosen uniformly over the localized null $\hat \Pi\in {\mP}_{n,k_0}$ 
In practice, we approximate the conditional quantiles by conditional Monte Carlo as described in Section \ref{subsubsection:testing permutation}.




\subsection{A finite-sample valid confidence set for the coefficient vector}

We next state the finite-sample coverage guarantee for the confidence region
constructed from the candidate permutation set $\mathcal{C}^{(L)}$.

\begin{theorem}[Coverage probability of the confidence set]\label{thm:confidence}
Suppose $C_{\min}(\Pi_0) > 0$ and $n > p + 1$ hold. Then, it holds that
\begin{equation}
         \bP_{\left(\mathcal{U}^L, \bm{U}\right)} \left( \bm{\beta}_0 \in \Gamma_{\alpha}^{\mathrm{coef}}(\bm{Y}^{\obs}) \right) \ge \alpha - \delta_L,
\end{equation}
where $\delta_L$ is specified in Theorem \ref{thm:recovery}.
\end{theorem}

Theorem~\ref{thm:confidence} establishes that the coverage of $\Gamma_{\alpha}^{\mathrm{coef}}(\bm{Y}^{\obs})$ for $\bm{\beta}_0$ up to two small, vanishing errors: one from the localization for the candidate set and another from Monte Carlo approximation.  
As $L$ grows, the coverage converges to the nominal level~$\alpha$.  

Theorem \ref{thm:confidence} represents a \emph{locally simultaneous} form of selective inference, which is conducted uniformly over a localized, data-dependent neighborhood of plausible permutations $\hat{\mP}_{n,k_0}$ that contains the true alignment with high probability (as shown in Theorem \ref{thm:recovery}), instead of conditioning on a single realized permutation. Within this neighborhood, classical reference laws such as the $F$ distribution remain valid, providing finite-sample coverage without the need to describe the full combinatorial selection event. At the same time, it highlights a notion of robustness to alignment uncertainty. In particular, if the true data-generating alignment is not uniquely identifiable but lies within this candidate set (as in record linkage, partial mismatching, or anonymization), then inference on $\beta_0$ still remains valid.


\subsection{Algorithmic guarantees}

In this section, we show that the estimator defined as the solution to the LAP \eqref{def:tilde_Pi} proposed in Section \ref{sec:algorithm} 
coincides with the solution to the original least square problem \eqref{eq:repro_Pi} with high probability. 
To clarify the dependence of the solutions on $\bm{U}$ and $\bm{U}^*$, we denote the solution to the LAP \eqref{def:tilde_Pi} as $\tilde{\Pi}(\bm{U}, \bm{U}^*)$, and the solution to the original problem \eqref{eq:repro_Pi} as $\hat{\Pi}(\bm{U}, \bm{U}^*)$.



\begin{theorem}[High-probability equivalence]\label{thm:high-probability equivalence}
For $\xi \in (0,1)$, let $B_Y(\xi), B_{\rm diag}(\xi), \eta_{\rm op}(\xi), \underline\Delta_\xi$ be the positive constants defined later in Section~\ref{subsec: proof of hungarian}.
Choose $\lambda_1,\lambda_2\ge 0$ such that
\begin{equation}\label{eq:lambda-window}
\lambda_1+\lambda_2 B_{\rm diag}(\xi) + \eta_{\rm op}(\xi)
 < \frac{n \underline\Delta_\xi}{2k}
\end{equation}
with the condition $\lambda_1\ge B_Y(\xi)/k$.
Then, we have
\begin{equation}
    \bP_{\bm{U}, \bm{U}^*}\left(\hat{\Pi}(\bm{U}, \bm{U}^*) = \tilde{\Pi}(\bm{U}, \bm{U}^*) \right) \ge 1 - 4\xi.
\end{equation}
\end{theorem}


About the constraint \eqref{eq:lambda-window}, $B_{\rm diag}(\xi)$ is the upper bound of the largest diagonal residual and $\eta_{\rm op}(\xi)$
controls the projection-mismatch error (see \eqref{eq:bdiag} and \eqref{eq:eta-op-def}),
while $\underline\Delta_\xi$ is an explicit lower bound on the gap of the objective function
(See Proposition \ref{prop:F separation gap}). 


Theorem~\ref{thm:high-probability equivalence} establishes a nonasymptotic argmin equivalence between a single linear assignment problem and the least-squares objective. Specifically, under an explicit finite-sample margin condition (quantified by $\underline{\Delta}_\xi$) and a transparent tuning window \eqref{eq:lambda-window}, the solution of the score-weighted LAP is identical to the global minimizer of the least square problem with probability at least $1-4\xi$.
The result provides an exact, rather than approximate, optimization correspondence under explicitly verifiable finite-sample conditions.
To our knowledge, this is the first formal proof of finite-sample equivalence between an $O(n^3)$-time assignment solver (Hungarian algorithm) and the global minimizer of a nonlinear least-squares problem defined over the combinatorially large permutation space.


\section{Extensions}\label{sec:extensions}

\subsection{Partially permuted designs}\label{subsec:partial-design}

In some applications the covariates can be decomposed as $(\bm X,\bm Z)$, where only $\bm X$ is
misaligned with the response.  We consider the partially permuted
Gaussian linear regression model
\begin{equation}\label{eq:model_partial}
\bm Y = \Pi_0 \bm X \bm\beta_{1,0} + \bm Z \bm\beta_{2,0} + \sigma_0 \bm U,
\qquad
\bm U\sim N(0,I_n),
\end{equation}
where $\bm X\in\mathbb R^{n\times p_1}$ and $\bm Z\in\mathbb R^{n\times p_2}$ are fixed.

For any $\Pi\in\mP_{n,k}$, define the expanded design matrix
\begin{equation}\label{eq:Wpartial-def}
W(\Pi):=\big[\Pi\bm X,\ \bm Z\big]\in\mathbb R^{n\times (p_1+p_2)},
\qquad
M_{W(\Pi)}:=W(\Pi)\left(W(\Pi)^\top W(\Pi)\right)^{-1}W(\Pi)^\top.
\end{equation}
Then the objective \eqref{eq:repro_Pi} becomes
\begin{equation}\label{eq:repro_Pi_partial}
F_{\rm partial}(\Pi):=\left\| (I_n-M_{W(\Pi)})\bm Y^{\obs}\right\|_2^2.
\end{equation}

Permutation candidates are generated exactly as in Section \ref{subsec:candidate}
by augmenting $W(\Pi)$ with the random direction $\bm u_\ell^*$:
\[
\widehat\Pi_{\ell}^{\rm partial}
\in \arg\min_{\Pi\in\mP_{n,k}}
\left\| (I_n - M_{(W(\Pi),\bm u_{\ell}^*)})\bm Y^{\obs}\right\|_2^2,
\qquad
\ell=1,\dots,L,
\]
and we define the candidate set
\begin{equation}\label{eq:perm-candset-partial}
\mathcal{C}_{\rm partial}^{(L)}
:=\left\{
\Pi\in\mP_{n,k}:\ \exists \ell\in[L]\ \text{s.t.}\ \widehat\Pi_{\ell}^{\rm partial}=\Pi
\right\}.
\end{equation}
Theorem~\ref{thm:recovery partial} shows that, under a
separation condition $C_{\min}^{\rm partial}(\Pi_0)>0$, the miscoverage probability
$\bP(\Pi_0\notin \mathcal{C}_{\rm partial}^{(L)})$ is bounded by
\[
\delta_L^{\rm partial}
=O\left(
\frac{(\log L)^{(n-(p_1+p_2))/2}}{L^{((n-(p_1+p_2))-1)/(2(n-1))}}
\right).
\]

Given $\Pi\in\mathcal{C}_{\rm partial}^{(L)}$, a joint confidence set for $(\bm\beta_1,\bm\beta_2)$ 
can be built using the standard $F_{p_1+p_2,n-(p_1+p_2)}$-distribution based on $W(\Pi)$:
\begin{equation}\label{eq:Tjointpartial-def}
T(\bm Y^{\obs},(\Pi,\bm\beta_1,\bm\beta_2))
:=\frac{n-(p_1+p_2)}{p_1+p_2}
\frac{\left(\bm Y^{\obs}-W(\Pi)\binom{\bm\beta_1}{\bm\beta_2}\right)^\top M_{W(\Pi)}
      \left(\bm Y^{\obs}-W(\Pi)\binom{\bm\beta_1}{\bm\beta_2}\right)}
     {\left(\bm Y^{\obs}-W(\Pi)\binom{\bm\beta_1}{\bm\beta_2}\right)^\top (I_n-M_{W(\Pi)})
      \left(\bm Y^{\obs}-W(\Pi)\binom{\bm\beta_1}{\bm\beta_2}\right)},
\end{equation}
and resulting $\alpha$-level confidence set is given by
\begin{equation}\label{eq:confset-partial-joint}
\Gamma_{\alpha}^{\mathrm{joint}}(\bm{Y}^{\obs})
:=\bigcup_{\Pi \in \mathcal{C}_{\mathrm{partial}}^{(L)}} \left\{
(\bm\beta,\bm\beta')\in \R^{p_1+p_2}:\ 
T(\bm Y^{\obs},(\Pi,\bm\beta,\bm\beta')) \le F_{p_1+p_2,n-(p_1+p_2)}^{-1}(\alpha)
\right\}.
\end{equation}

If the main interest is $\bm\beta_1$ with $\bm\beta_2$ treated as a nuisance parameter, we use the
partialled--out $F$-distribution.  Let $M_{\bm Z}$ be the orthogonal projector onto $\range(\bm Z)$ and
set $P_{\bm Z}:=I_n-M_{\bm Z}$. For any $\Pi\in\mP_{n,k}$ and $\bm\beta_1\in\mathbb R^{p_1}$,
define
\begin{equation}\label{eq:Tpartial-def}
T(\bm Y^{\obs},(\Pi,\bm\beta_1)\mid \bm Z)
:=\frac{n-(p_1+p_2)}{p_1}
\frac{
(\bm Y^{\obs}-\Pi\bm X\bm\beta_1)^\top (M_{W(\Pi)}-M_{\bm Z}) (\bm Y^{\obs}-\Pi\bm X\bm\beta_1)
}{
(\bm Y^{\obs}-\Pi\bm X\bm\beta_1)^\top (I_n-M_{W(\Pi)}) (\bm Y^{\obs}-\Pi\bm X\bm\beta_1)
}.
\end{equation}
The resulting $\alpha$-level confidence set is given by
\[
\Gamma_{\alpha}^{\rm partial}(\bm Y^{\obs})
:=\bigcup_{\Pi\in\mathcal C_{\rm partial}^{(L)}}
\left\{
\bm\beta_1\in\mathbb R^{p_1}:\ 
T(\bm Y^{\obs},(\Pi,\bm\beta_1)\mid \bm Z)\le F_{p_1,n-(p_1+p_2)}^{-1}(\alpha)
\right\},
\]
and its coverage guarantee is shown in Theorem~\ref{thm:confidence set partial beta1}.

\subsection{Permuted ridge regression}
Our methodology can also be extended to a permuted regression model with a ridge penalty.
Consider the data generating process analogous to~\eqref{eq:random},
\[
\bm Y^{\mathrm{obs}}=\Pi_{0}\bm X\bm{\beta}_{0}+\sigma_{0}\bm U,
\qquad \Pi_0\in\mP_{n,k}.
\]
For a fixed $\lambda>0$, define the ridge objective
\[
F_\lambda(\Pi):=\min_{\bm{\beta}\in\R^p}
\left\{\|\bm Y^{\mathrm{obs}}-\Pi\bm X\bm{\beta}\|_2^2+\lambda\|\bm{\beta}\|_2^2\right\},
\qquad \Pi\in\mP_{n,k}.
\]

It is convenient to rewrite the ridge objective in an augmented least-squares form
(e.g.,~\cite{hastie2020ridge}). Define
\[
\bm Y^{\mathrm{aug}}:=\begin{pmatrix}\bm Y^{\mathrm{obs}}\\ \bm 0_p\end{pmatrix}\in\R^{n+p},
\qquad
\bm X^{\mathrm{aug}}:=\begin{pmatrix}\bm X\\ \sqrt{\lambda}I_p\end{pmatrix}\in\R^{(n+p)\times p},
\qquad
\tilde\Pi:=\begin{pmatrix}\Pi & 0\\ 0 & I_p\end{pmatrix}\in\mathcal P_{n+p}.
\]
Then $F_\lambda(\Pi)=\min_{\bm{\beta}}\|\bm Y^{\mathrm{aug}}-\tilde\Pi\bm X^{\mathrm{aug}}\bm{\beta}\|_2^2$.
Let $\tilde\Pi_0$ denote the block-diagonal embedding of $\Pi_0$, namely
\[
\tilde\Pi_0 := \begin{pmatrix}\Pi_0 & 0\\ 0 & I_p\end{pmatrix}\in\mathcal P_{n+p},
\]
and define the augmented noise $\tilde{\bm U}:=(\bm U,\bm 0_p)\in\R^{n+p}$.
Under the model we have $\bm Y^{\mathrm{aug}}=\tilde\Pi_0\bm X^{\mathrm{aug}}\bm{\beta}_0+\sigma_0\tilde{\bm U}$.
For $\Pi\in\mP_{n,k}$ write $
W_\lambda(\Pi):=\tilde\Pi\bm X^{\mathrm{aug}}\in\R^{(n+p)\times p}$.

For i.i.d.\ $\bm U^*_1,\dots,\bm U^*_L\sim N(0,I_n)$ let
$\tilde{\bm U}^*_\ell:=(\bm U^*_\ell,\bm 0_p)$ and define
\begin{equation}\label{eq:perm_candidates_ridge}
\hat\Pi_{\lambda,\ell}\in
\arg\min_{\Pi\in\mP_{n,k}}
\left\|
(I_{n+p}-M_{(W_\lambda(\Pi),\tilde{\bm U}^*_\ell)})\bm Y^{\mathrm{aug}}
\right\|_2^2,
\qquad
\mC_\lambda^{(L)}:=\{\hat\Pi_{\lambda,1},\dots,\hat\Pi_{\lambda,L}\}.
\end{equation}

\begin{remark} In the augmented representation, the last $p$ pseudo-observations are noiseless, so the projected augmented noise is no longer exactly spherically symmetric when $\lambda$ is very small. To keep the theory parallel to the linear regression case, we impose the mild non-degeneracy condition \begin{equation}\label{eq:ridge_lambda_lower} c_\lambda:=\frac{\lambda}{\lambda+\|\bm X\|_{\mathrm{op}}^2}\ge \underline c\in(0,1), \end{equation} under which the argument used in Lemma~\ref{lem:max-sim} carries over up to constants (see Lemma~\ref{lem:max_sim_ridge} in Appendix~\ref{subsec:ridge_theory}). \end{remark}


\section{Selection of tuning parameters}\label{sec:tuning}

We propose a selection method of the two regularization parameters $(\lambda_1,\lambda_2)$ for the score-weighted LAP in Section~\ref{sec:algorithm}. 

Specifically, we select the parameters as satisfying the constraint \eqref{eq:lambda-window} in Theorem \ref{thm:high-probability equivalence}. 
Throughout this section, we fix $\xi=0.05$.

The constraint \eqref{eq:lambda-window} is stated in terms of the explicit lower bound
$n\underline\Delta_\xi$.
At the proof level, the comparison is with the (random) separation gap
$\Delta_F$, which will be defined in \eqref{eq:delta-F}, and Proposition~\ref{prop:F separation gap} shows that
$\Delta_F\ge n\underline\Delta_\xi$ with probability at least $1-\xi$.
Motivated by this connection, our implementation chooses $(\lambda_1,\lambda_2)$ via simple
plug-in estimates of $B_{\rm diag}(\xi)$, $\eta_{\rm op}(\xi)$, and $\Delta_F$.

\begin{enumerate}
\item Compute the empirical analogs $\widehat{\bm m} = M_{(\bm X,\bm u^*)}\bm Y^{\obs}$ and set $
      \widehat B_{\rm diag}:=\max_{1\le i\le n}\left(Y_i^{\obs}-\widehat m_i\right)^2$.
\item Evaluate the closed-form expression of $\eta_{\rm op}(\xi)$ in \eqref{eq:eta-op-def} by
      replacing the unknown norms by the observable quantity $\|\bm Y^{\obs}\|_2^2$,
      yielding $\widehat\eta_{\rm op}$.
\item Estimate the separation gap $\Delta_F$ in \eqref{eq:delta-F} as follows.
      First, compute a baseline $k$-sparse permutation $\Pi_{\rm base}$ as a solution of the score-weighted LAP with $(\lambda_1,\lambda_2)=(0,0)$. Next, generate a small collection of
      neighbouring permutations by applying random swaps to two rows of $\Pi_{\rm base}$, and compute
      $F(\Pi)-F(\Pi_{\rm base})$ for these candidates. We set $\widehat\Delta_F$ to be the minimum of
      the positive differences.
    
\item Define an available budget $
      b:=0.9(\frac{\widehat\Delta_F}{2k}-\widehat\eta_{\rm op})_+$,
      where the constant $0.9$ is a fixed safety factor to hedge against estimation error.
      Finally, split the budget equally between $\lambda_1$ and $\lambda_2$ as $
      \lambda_1=\frac{b}{2}$ and $\lambda_2=\frac{b}{2\max\{\widehat B_{\rm diag},10^{-12}\}}$.
\end{enumerate}

Through our experiments in Section \ref{sec:numerical}, we adopt these tuning parameters as a default rule. 



\section{Numerical results}\label{sec:numerical}

\subsection{Simulation analysis}
In this section, we assess the finite-sample performance of the proposed method.
We focus on three aspects:
(i) the size and coverage probability of the candidate set (Section \ref{subsec:candidate});
(ii) the size and power of the conditional test for the existence of a permutation (Section \ref{subsubsection:testing permutation}); and
(iii) the coverage probability and the volume of the confidence set for the coefficient vector (Section \ref{subsec:coef}).

All experiments are based on the shuffled linear regression model
\[
\bm{Y} = \Pi_0 \bm{X}\bm{\beta}_0 + \sigma_0 \bm{U}, \qquad \bm{U}\sim\Normal(\bm{0}, I_n),
\]
where $\bm{\beta}_0=(0.5,-1,2)^\top$ and the design matrix $\bm{X} \sim\Normal(\bm{0},I_p)$ with $p=3$.
For each configuration, $\Pi_0$ and $\bm{X}$ are fixed across Monte Carlo replications.

We vary the sample size $n\in\{100,200,400\}$, the number of mismatches $k\in\{0,2,5,10,15,20\}$,
the noise level $\sigma_0\in\{0.01,0.05,0.10,0.20,0.30,0.40,0.50\}$,
and the number of repro samples $L\in\{200,400\}$.
Each setting is replicated $200$ times.
The repro step employs the score-weighted LAP (Algorithm~\ref{alg:scoring hungarian})
with tuning parameters $(\xi,\rho,\alpha_{\mathrm{diag}})=(0.05,0.9,0.5)$ specified in Section \ref{sec:tuning}.

\subsubsection{Performance of the candidate set}
We assess the empirical performance of the candidate set $\mathcal{C}^{(L)}$.
For each configuration, we compute the average size of $\mathcal{C}^{(L)}$ (the number of unique permutations) across Monte Carlo replications, the empirical probability that the true permutation $\Pi_0$ is contained in $\mathcal{C}^{(L)}$, and the average matching fraction defined below.

As summarized in Figures~\ref{fig:size candidate},
the size of $\mathcal{C}^{(L)}$  increases with the noise level $\sigma_0$ and the mismatch rate $k$,
reflecting greater combinatorial ambiguity under weaker signal or stronger misalignment, in contrast, larger $n$ yields smaller and more stable candidate sets.
To evaluate the alignment quality, we report a matching-fraction diagnostic that measures
how closely the best candidate permutation aligns with the true $\Pi_0$.
For each Monte Carlo iteration, we compute the normalized complement of the
Hamming distance $
\psi = 1 - \frac{1}{n} \min_{1 \le \ell \le L} d(\widehat{\Pi}_\ell, \Pi_0)$,
which represents the proportion of correctly matched indices in the best case among the $L$ repro samples.
We then average $\psi$ across the Monte Carlo replications and report it in Figure~\ref{fig:fraction}.
The matching fraction decreases with increasing $\sigma_0$ and $k$,
improves with larger $n$, and approaches $100\%$ when $k=0$.
Although perfect recovery is generally hard due to the discrete nature of permutation matrices,
the candidate sets achieve sufficiently high matching fractions,
indicating that most columns are correctly aligned even under moderate noise.
These results confirm that repro-sample localization effectively narrows the permutation space
while maintaining near-perfect inclusion of the true $\Pi_0$.



\subsubsection{Size and power of the test for the existence of a permutation}
Next, we assess the finite-sample size and power of the conditional Monte Carlo test described in Section \ref{subsubsection:testing permutation}.
We use the same baseline model and parameter grid as above.
For each configuration, we record whether the null $H_0:\Pi_0=I_n$ is rejected at nominal level $\alpha=0.05$.
The empirical rejection rates, averaged over $200$ repetitions,
are plotted as a function of $\sigma_0$ in Appendix~Figure~\ref{fig:test size power}.
When $k=0$ (no mismatch), the rejection frequency remains at or below the nominal significance level across all $\sigma_0$,
demonstrating valid Type-I error control guaranteed by Theorem~\ref{thm:p-value}.
As the number of mismatches $k$ increases, the rejection probability rises monotonically, indicating increasing power under stronger alternatives.
Larger sample sizes $n$ yield slightly higher power at moderate noise levels, reflecting the greater information available for detecting mismatches.
Differences between $L=200$ and $L=400$ are modest.


\subsubsection{Coverage probability of the confidence set for the coefficient vector}
Finally, we evaluate the finite-sample coverage and size of the confidence sets
constructed for the regression coefficient vector $\bm{\beta}_0$
using the union method described in Section~\ref{subsec:coef}.
For each configuration, we construct the confidence set $
\Gamma_\alpha^{\mathrm{coef}}(\bm{Y}^{\mathrm{obs}})$ with $\alpha=0.95$, and record whether $\bm{\beta}_0$ lies inside the region as well as its numerical volume.

The results in Figures \ref{fig:coverage coef} shows that the empirical coverage probability closely matches the nominal level across all noise regimes,
verifying the finite-sample validity guaranteed by Theorem~\ref{thm:confidence}. 
In addition to the scalar summaries of coverage, we visualize the geometry of the constructed confidence region.
Figure~\ref{fig:repro_union_three_planes} shows three two-dimensional projections of the union confidence set
onto the coordinate planes
for one representative configuration ($n=100$, $k=5$, $L=200$, $\sigma_0=0.10$).
In all projections, the true parameter $\bm{\beta}_0=(0.5,-1,2)^\top$ lies well within the confidence region,
demonstrating that the procedure is consistent with the finite-sample guarantee.
Moreover, the projected regions are compact and tightly concentrated around the true value,
illustrating that taking unions over $\mC^{(L)}$
does not lead to excessive conservatism: the volume of confidence regions are also reported in Figure \ref{fig:app-coef-volume}.

\begin{figure}[htbp]
    \centering
    \begin{subfigure}[b]{0.32\textwidth}
        \centering
        \includegraphics[width=\textwidth]{figures/sim_sigma01_n100_k5_L200_union_true_01.png}
        \caption{Projection onto $(\beta_1,\ \beta_2)$}
    \end{subfigure}
    \hfill
    \begin{subfigure}[b]{0.32\textwidth}
        \centering
        \includegraphics[width=\textwidth]{figures/sim_sigma01_n100_k5_L200_union_true_12.png}
        \caption{Projection onto $(\beta_2,\ \beta_3)$}
    \end{subfigure}
    \hfill
    \begin{subfigure}[b]{0.32\textwidth}
        \centering
        \includegraphics[width=\textwidth]{figures/sim_sigma01_n100_k5_L200_union_true_02.png}
        \caption{Projection onto $(\beta_1,\ \beta_3)$}
    \end{subfigure}
    \caption{Two-dimensional projections of the union confidence set.}
    \label{fig:repro_union_three_planes}
\end{figure}

\begin{figure}[htbp]
  \centering
  % --- (a)
  \begin{subfigure}[t]{0.48\linewidth}
    \centering
    \includegraphics[width=\linewidth]{figures/mc_unique_Pi_count_all.png}
    \caption{Unique permutations in $\mC^{(L)}$.}
    \label{fig:size candidate}
  \end{subfigure}\hfill
  % --- (b)
  \begin{subfigure}[t]{0.48\linewidth}
    \centering
        \includegraphics[width=\linewidth]{figures/match_fraction_best_all.png}
    \caption{Matching fraction (\%).}
    \label{fig:fraction}
  \end{subfigure}

  \vspace{0.6em}

  % --- (c)
  \begin{subfigure}[t]{0.48\linewidth}
    \centering
    \includegraphics[width=\linewidth]{figures/test_reject_all.png}
    \caption{Rejection rate for $k=0$ (size) and power for $k>0$; nominal $\alpha=0.05$ shown (red dotted).}
    \label{fig:test size power}
  \end{subfigure}\hfill
  % --- (d)
  \begin{subfigure}[t]{0.48\linewidth}
    \centering
    \includegraphics[width=\linewidth]{figures/coef_covered_all.png}
    \caption{Coverage probability of coefficient regions; nominal $\alpha=0.95$ shown (red dotted).}
    \label{fig:coverage coef}
  \end{subfigure}
  \caption{Characteristics of the candidate set $\mC^{(L)}$.  \label{fig:simulations}}
\end{figure}




\subsection{Empirical application: Beijing air-quality data}\label{subsec:real data}

We illustrate the proposed finite-sample inference for shuffled regression using the Beijing Multi-Site Air Quality dataset \citep{beijing_multi-site_air_quality_501}. The objective is to verify that the permutation test behaves consistently with the theoretical guarantees. To this end, we construct two controlled scenarios on the real data: a baseline with no mismatch and a perturbed setting in which we artificially shuffle the response–covariate matching within short time windows. The dataset reports hourly concentrations of six major air pollutants together with meteorological covariates from twelve stations between March~2013 and February~2017.  
Following the empirical design of \cite{slawski2020two}, we focus on the Nongzhanguan station and restrict the sample to the years 2016–2017.  
After removing rows with missing entries, we standardize each variable to zero mean and unit variance and extract a contiguous window of $n=1000$ hourly observations (approximately six weeks of data).  
The response variable $\bm{Y}$ is the concentration of PM$_{2.5}$, and the design matrix $\bm X$ contains $p=10$ continuous covariates: PM$_{10}$, SO$_2$, NO$_2$, CO, O$_3$, temperature, dew point, pressure, precipitation, and wind speed.  
All computations employ the score-weighted LAP described in Section~\ref{sec:methodology} within the repro samples framework.

We analyze two complementary scenarios.  
First, a sparse and local mismatch is introduced: $8\%$ of the rows of $\bm Y$ are permuted within the same day and within a $\pm 3$-hour window to emulate record-linkage or time-synchronization errors, as in \cite{slawski2020two}.  
Second, no mismatch is added, so that the alignment between $\bm Y$ and $\bm X$ remains exact.  For both settings, the candidate set $\mC^{(L)}$ is built from $L=250$ repro samples, and the conditional Monte Carlo test of $H_0:\Pi_0=I_n$ is conducted at nominal level $\alpha = 0.05$. 

Under the local-mismatch configuration, the candidate set expands to $|\mC^{(L)}|=224$ and $H_0$ is rejected at $\alpha = 0.05$, indicating statistically detectable misalignment.  
Under the no-mismatch configuration, the procedure produces a singleton candidate set $\mC^{(L)}=\{I_n\}$ and the null $H_0$ is not rejected, confirming 
Type-I error control in accordance with Theorem~\ref{thm:p-value}.  

These results are consistent with the theoretical predictions: when the data contain no mismatches, the candidate set collapses to the identity, whereas local shuffling enlarges the set and triggers rejection through the conditional test.  

Overall, the two real-data scenarios demonstrate that the proposed finite-sample inference is fully aligned with the theoretical properties established for the repro samples framework.



\begin{table}[htbp]
\centering
\caption{Results for the Beijing air-quality dataset (Nongzhanguan St., 2016–2017).}
\label{tab:realdata}
\begin{tabular}{lccc}
\toprule
\textbf{Scenario} & \textbf{Mismatch rate} & \textbf{Candidate set size} $|\mC^{(L)}|$ & \textbf{Test result} \\
\midrule
Local mismatch & 8\% (within $\pm 3$h) & 224 & Rejected ($p<0.05$) \\
No mismatch & 0\% & 1 ($\{I_n\}$) & Not rejected ($p>0.05$) \\

\bottomrule
\end{tabular}
\end{table}





\section{Conclusion}\label{sec:conclusion}
This paper develops a finite-sample inference framework for linear regression with sparsely permuted data. Building upon the repro sample method and a localization step, the method converts an intractable search over $\mathcal{P}_{n,k}$ into inference that is uniform over a small, data-dependent candidate set $\mathcal{C}^{(L)}$.

Our theory shows: (i) the miscoverage probability of the candidate set $\bP\{\Pi_0\notin \mathcal{C}^{(L)}\}$ decays polynomially in $L$ as the number of repro samples grows (Theorem~\ref{thm:recovery}); (ii) a conditional Monte Carlo test for the $k_0$-sparsity with finite-sample type-I error control (Theorem~\ref{thm:p-value}); (iii) confidence regions for $\bm \beta_0$ obtained by taking union over the candidate set $\mathcal{C}^{(L)}$ with finite-sample coverage (Theorem~\ref{thm:confidence}); and (iv) a computational guarantee that the proposed score-weighted LAP attains the same global minimizer as the least square problem with high probability (Theorem~\ref{thm:high-probability equivalence}) with suitable tuning parameters. Extensions to partially permuted designs and ridge regularization are also discussed.

Empirically, a few hundred repro samples deliver near-nominal behavior: the test controls the type-I error, the union region is robust to mild misalignment, and computation scales as $O(n^3)$ with tuning that ensures the argmin equivalence.

Future research includes two main directions. First, extending the calibration beyond Gaussian errors to sub-Gaussian or heavy-tailed settings would enhance robustness. Second, exploring high-dimensional extensions with structured regularization and multi-response \cite{slawski2020two} or block-permutation designs \cite{abbasi2025alternating} would broaden applicability. Overall, the results demonstrate that locally simultaneous, repro samples-based inference offers a general and computationally feasible foundation for exact finite-sample validity in irregular, combinatorial statistical models including permutations.

\section*{Acknowledgements} 
The authors are deeply grateful to Shinji Ito, Kaoru Irie and Takeru Matsuda for their helpful discussions and insightful comments, which greatly improved the clarity and scope of this work. We also thank all members of our research group for their constructive feedback during the early stages of this project. Hirofumi Ota is supported by JSPS KAKENHI (25K23133). Masaaki Imaizumi is supported by JSPS KAKENHI (24K02904), JST CREST (JPMJCR21D2),  JST FOREST (JPMJFR216I), and JST BOOST (JPMJBY24A9).


\newpage
\appendix

\section{Additional Methodological Background}

\subsection{Repro samples method}
In this subsection, we give a brief review on the repro samples framework to obtain a performance-guaranteed confidence set on general parameters. The \textit{repro samples method} \cite{xie2022repro, xie2024repro} is a simulation-inspired general inference method, which is particularly effective for making inference on irregular inference problems \cite{wasserman2020universal}. The general setup is as follows: suppose the sample data $\bm{Y} \in \mathcal{Y}$ are generated from a model with the true parameter $\bm{\theta}_{0} \in \Theta$. We also assume that we know how to simulate $\bm{Y}$ from the statistical model indexed by $\bm{\theta} \in \Theta$ and it has the form of generative algorithmic model: 
\begin{equation}\label{eq:algorithmic}
    \bm{Y} = G(\bm{\theta}, \bm{U}),
\end{equation}
where $\bm{U} = (U_1, \cdots, U_r) \in \mathcal{U} \subset \R^r$ is a random vector whose distribution is free of $\bm{\theta} \in \Theta$ and $G: \Theta \times \mathcal{U} \mapsto \mY$ is a known mapping. Let 
\begin{equation}\label{eq:realization}
    \bm{Y}^{\mathrm{obs}} = G(\bm{\theta}_{0}, \bm{u}^{\mathrm{rel}})
\end{equation}
be the realized version of $\bm{Y}$ and also $\bm{u}^{\mathrm{rel}}$ be the unobserved realization of $\bm{U}$.

The general framework of repro samples method to construct a performance-guaranteed confidence set is as follows. We generate $\bm{u}^{*} \sim \bm{U}$ given $\bm{\theta}^{*}$ based on the generative algorithmic model (\ref{eq:algorithmic}), and call $\bm{Y}^{*} = G(\bm{\theta}^{*}, \bm{u}^{*})$ as the repro samples. The key idea to construct a valid confidence set is to quantify the uncertainty of $\bm{U}$ through $\bm{u}^{*}$. Let $B_\alpha \subset \mathcal{U}$ be a level-$\alpha$ Borel set for some $\alpha \in (0,1)$, such that $\bP(\bm{U} \in B_\alpha) \ge \alpha$. Then, we define 
\begin{equation*}
    \Gamma_\alpha(\bm{Y}^{\mathrm{obs}}) = \left\{ \bm{\theta} : \exists \bm{u}^{*} \in B_\alpha \ \mathrm{s.t.}\ \bm{Y}^{\mathrm{obs}} = G(\bm{\theta}, \bm{u}^{*}) \right\}.
\end{equation*}
It is shown that the $\Gamma_\alpha(\bm{Y}^{\mathrm{obs}})$ is a level-$\alpha$ confidence set for $\bm{\theta}_{0}$, i.e.,
    \[
    \bP(\bm{\theta}_{0} \in \Gamma_\alpha(\bm{Y})) \ge \alpha,\ \hbox{for }\alpha \in (0,1).
    \]
In the simplest case above, we only need to quantify the uncertainty of $\bm{U}$ via a level-$\alpha$ Borel set $B_\alpha$. More generally, we can consider the case where a nuclear mapping $T: \mathcal{U} \times \Theta \mapsto \R^d$ is available. Let $B_\alpha(\bm{\theta})$ be a level-$\alpha$ Borel set such that $\bP(T(\bm{U}, \bm{\theta}) \in B_\alpha(\bm{\theta}) )\ge \alpha$. Then, the following
\begin{equation}\label{eq:confidence set2}
    \tilde{\Gamma}_\alpha(\bm{Y}^{\mathrm{obs}}) = \left\{ \bm{\theta} : \exists \bm{u}^{*} \ \mathrm{s.t.}\ \bm{Y}^{\mathrm{obs}} = G(\bm{\theta}, \bm{u}^{*}),  T(\bm{u}^*, \bm{\theta}) \in B_\alpha(\bm{\theta})\right\}
\end{equation}
is also a valid level-$\alpha$ confidence set, i.e., it holds that, 
\begin{equation}\label{eq:confidence set2 validity}
    \bP(\bm{\theta}_{0} \in  \tilde{\Gamma}_\alpha(\bm{Y})) \ge \bP(T(\bm{U}, \bm{\theta}) \in B_\alpha(\bm{\theta}) )\ge \alpha,\ \hbox{for }\alpha \in (0,1).
\end{equation}

\section{Connection to selective/post-selection inference}

The classical selective inference provides valid post–model‐selection inference by
conditioning on the event that a data–driven selection procedure chooses a particular model. For example, \cite{fithian2014optimal, lee2016exact, liu2018more} proceed statistical inference by conditioning on the realised event and deriving the exact conditional distribution of the test statistic under this restriction. The
conditional distribution is characterized as a truncated multivariate normal whose
support is a polyhedron, via the Gaussian polyhedral lemma \cite{lee2016exact}.
While the resulting procedures control coverage and type-I error conditional on the selected event, three practical difficulties are common: (i) the selection event is described by high-dimensional linear
      inequalities, so explicit formulas are rare and simulation is costly; (ii) conditional guarantees do not automatically imply unconditional
      guarantees so that aggressive selection can produce undercoverage when the conditioning is removed; (iii) conditioning discards information that would otherwise contribute to
      statistical power, especially in small samples.
      Another notable methodology for the selective inference is locally simultaneous inference \cite{zrnic2024locally}, 
which addresses these issues by replacing conditioning with simultaneous control over a data-dependent neighbourhood. Compared to the standard solution such as the Bonferroni-type method taking the unions across all candidate parameters, the locally simultaneous inference provides a valid uncertainty quantification over the data-depoendent localized parameters and therefore is less conservative.

While the conditional inference fixes the selected model and conditions on that event, the locally simultaneous inference keeps the selection random and enforces validity uniformly over the random set. Furthermore, the conditional methods theoretically and practically require a precise characterization of the conditional law which is often achieved by the Gaussian polyhedral lemma, while locally simultaneous methods use standard parametric (or possibly nonparametric) reference distributions and avoid any geometric description of the selection event. 

From these perspectives, we adopt the locally simultaneous inference for shuffled/permuted linear regression models. Our approach first provides the valid localization scheme based on the repro samples method: Theorem \ref{thm:recovery} establishes the true permutation is included in the candidate set of permutations with probability. Within this candidate set, we test the null hypothesis for the underlying permutation structure and construct a union confidence set for the regression coefficient. Our results demonstrate that locally simultaneous inference extends naturally to problems with discrete and combinatorial parameters.
The methodology is conceptually simple, requires only classical distribution
theory, and exhibits strong finite-sample performance, thereby providing a
practical alternative to conditional selective inference.

\section{Details of Hungarian Algorithm}\label{sec:hungarian alg}

Given a cost matrix $\Omega\in\R^{n\times n}$, the linear assignment problem (LAP) is given by
\begin{equation}\label{eq:LAP-appendix}
  \hat{\Pi}\ \in\ \argmin_{\Pi\in\mP_n} \langle\Omega,\Pi\rangle
  =\argmin_{\Pi\in\mP_n}\sum_{i=1}^n\sum_{j=1}^n \Omega_{ij}\Pi_{ij}.
\end{equation}
In Algorithm~\ref{alg:scoring hungarian}, we apply this LAP with the score-weighted cost matrix
$\Omega=\Omega(\bm u^*)$.

A convenient viewpoint is to consider the linear-programming relaxation
\[
\min_{\Pi\ge 0}\left\{\langle\Omega,\Pi\rangle:
\sum_{j=1}^n\Pi_{ij}=1\ \forall i,\ \sum_{i=1}^n\Pi_{ij}=1\ \forall j\right\},
\]
whose extreme points are exactly the permutation matrices. Hence the relaxation attains an optimizer in
$\mP_n$. Also, the dual problem is given by
\[
\max_{\bm a,\bm b\in\R^n}\left\{\sum_{i=1}^n a_i+\sum_{j=1}^n b_j:
a_i+b_j\le\Omega_{ij}\ \forall i,j\right\}.
\]

The Hungarian algorithm \cite{kuhn1955hungarian} solves \eqref{eq:LAP-appendix} by maintaining
(i) a dual-feasible pair $(\bm a,\bm b)$ with $a_i+b_j\le \Omega_{ij}$ for all $i,j$ and
(ii) a partial matching $\mathcal M\subseteq[n]\times[n]$ using only the \emph{tight} (equality) edges
\[
E(\bm a,\bm b) := \left\{(i,j)\in[n]\times[n]: a_i+b_j=\Omega_{ij}\right\}.
\]
For any set of rows $S\subseteq[n]$, define its tight-neighborhood
\[
N(S)\ :=\ \left\{j\in[n]:\ \exists i\in S\ \text{such that}\ (i,j)\in E(\bm a,\bm b)\right\}.
\]
At each iteration the algorithm either augments the matching along tight edges
or updates the dual variables to create at least one new tight edge, while preserving dual feasibility.
When a perfect matching on $E(\bm a,\bm b)$ is obtained, complementary slackness implies primal optimality.

\begin{algorithm}[H]
\caption{Hungarian algorithm for $\min_{\Pi\in\mP_n}\langle \Omega,\Pi\rangle$}
\label{alg:scoring hungarian}
\begin{algorithmic}[1]
\REQUIRE $\Omega\in\R^{n\times n}$
\ENSURE $\hat{\Pi}\in\mP_n$
\STATE $a_i\leftarrow \min_{j\in[n]} \Omega_{ij}$ for all $i\in[n]$; \quad
       $b_j\leftarrow 0$ for all $j\in[n]$; \quad $\mathcal M\leftarrow\emptyset$.
\WHILE{$|\mathcal M|<n$}
  \STATE Pick an unmatched row $i_0$; set $S\leftarrow\{i_0\}$; $T\leftarrow N(S)$.
  \WHILE{every $j\in T$ is matched in $\mathcal M$}
     \STATE $S\leftarrow S\cup\{i:\ (i,j)\in\mathcal M\ \text{for some }j\in T\}$;
            \quad $T\leftarrow N(S)$.
  \ENDWHILE
  \IF{there exists an unmatched column $j_0\in T$}
     \STATE Augment $\mathcal M$ along tight edges to match $i_0$ with $j_0$.
  \ELSE
     \STATE $\delta\leftarrow \min\{\Omega_{ij}-a_i-b_j:\ i\in S,\ j\notin T\}$.
     \STATE $a_i\leftarrow a_i+\delta$ for all $i\in S$; \quad $b_j\leftarrow b_j-\delta$ for all $j\in T$.
  \ENDIF
\ENDWHILE
\STATE Output $\hat{\Pi}$ induced by the perfect matching $\mathcal M$.
\end{algorithmic}
\end{algorithm}

The potentials $(\bm a,\bm b)$ define reduced costs $\Omega_{ij}-(a_i+b_j)\ge 0$,
and the tight-edge set $E(\bm a,\bm b)$ collects the zero reduced-cost edges.
The algorithm attempts to increase the matching size using only tight edges.
If this is not possible for the current row set $S$, the slack update with
\[
\delta=\min\{\Omega_{ij}-a_i-b_j:\ i\in S,\ j\notin T\}
\]
preserves dual feasibility and creates at least one new tight edge, allowing the process to continue.
Here, ``augment'' means flipping matched/unmatched status along an alternating path of tight edges,
which increases $|\mathcal M|$ by one. After $n$ augmentations a perfect matching is obtained, and
complementary slackness implies that the induced permutation matrix $\hat{\Pi}$ is optimal for
\eqref{eq:LAP-appendix}.

\section{Proofs}\label{sec:proofs}
We first define a similarity measure $\phi:\R^n \times \R^n \to [0,1]$ between $\bm{v}_1, \bm{v}_2 \in \R^n$ as the squared cosine between $\bm{v}_1$ and $\bm{v}_2$:
\begin{equation}
    \phi(\bm{v}_1, \bm{v}_2) = \frac{(\bm{M}_{\bm{v}_1}\bm{v}_2 )^2}{\|\bm{v}_2 \|^2_2} = \frac{(\bm{v}_1^\top\bm{v}_2 )^2}{\|\bm{v}_1\|^2_2 \|\bm{v}_2 \|^2_2}.
\end{equation}
The cosine similarity $\phi(\bm{v}_1, \bm{v}_2)$ approaches one as the angle between $\bm{v}_1$ and $\bm{v}_2$ decreases.

\subsection{Supportive lemmas}
\begin{lemma}[Lemma 11 in \cite{wang2022finite}]\label{lem:similarity measure}
    Suppose $\bm{U}_1^*, \cdots, \bm{U}_L^*$ are i.i.d. copies of $\bm{U}^* \sim N(0, I_n)$. Then, for any $\gamma_2^2 \in (0,1)$, it holds that
    \begin{equation}
    \bP\left( \bigcap_{\ell = 1}^{L} \{ \phi(\bm{U}^*_\ell, \bm{U}) \le 1 - \gamma_2^2 \} \right) \le \left( 1- \frac{\gamma_2^{n-1}}{n-1} \right)^L. 
    \end{equation}
\end{lemma}
\begin{proof}
    See the proof of Lemma 11 in \cite{wang2022finite}.
\end{proof}
\begin{lemma}\label{lem:k-sparse permutation}
    The number of $k$-sparse permutation matrices $|\mP_{n,k}|$ is exactly calculated as follows:
    \begin{equation}
        |\mP_{n,k}| = \sum_{m = n - k}^{n} \binom{n}{m} D_{n - m},
    \end{equation}
    where $D_{n-m}$ is the number of permutations with no fixed points, computed as
    \[
D_{n-m} = (n-m)! \sum_{j=0}^{n-m} \frac{(-1)^j}{j!}.
\]
\end{lemma}
\begin{proof}
A permutation has exactly $m$ fixed points if we (i) choose the fixed-point set $T\subset[n]$ with $|T|=m$ (there are $\binom{n}{m}$ choices),
and (ii) derange the remaining $n-m$ indices (there are $D_{n-m}$ choices).
A $k$-sparse permutation satisfies $d(\Pi,I_n)\le k$, i.e., it has at least $n-k$ fixed points, so $m$ ranges from $n-k$ to $n$.
Summing over $m=n-k,\dots,n$ yields the formula.
\end{proof}

\begin{lemma}\label{lem:F-union}
Define $\Psi(x)=x-1-\log x$ for $x>0$. Fix $\gamma_1,\gamma_2\in(0,1)$ and define the event
\[
\mathcal E(\gamma_1,\gamma_2)
:=
\left\{
\max_{\Pi\in\mP_{n,k}}
\phi\left((I_n-M_{\Pi\bm X})\bm U^*, (I_n-M_{\Pi\bm X})\Pi_0\bm X\bm\beta_0\right)
<\gamma_1^2,\ 
\phi(\bm U^*,\bm U)>1-\gamma_2^2
\right\}.
\]
Assume the following holds:
\[
\min\left\{
\frac{1-\gamma_1^2}{2\sigma_0^2\gamma_2^2},\
\frac{(1-\gamma_1^2)^2}{16\sigma_0^2\gamma_2^2}
\right\}C_{\min}(\Pi_0)>1.
\]
Then, we have
\begin{align*}
&\bP_{(\bm U^*,\bm U)}\left(
\exists\Pi\in\mP_{n,k}\setminus\{\Pi_0\}:\ 
F(\Pi,\bm U^*)<F(\Pi_0,\bm U^*)
\mid 
\mathcal E(\gamma_1,\gamma_2)
\right)\\
&\le
|\mP_{n,k}|\left\{
\exp\left( -\frac{n}{2}\Psi\left( \frac{(1-\gamma_1^2)C_{\min}(\Pi_0)}{2\sigma_0^2\gamma_2^2}\right)\right)
+\exp\left( -\frac{n}{2}\Psi\left( \frac{(1-\gamma_1^2)^2C_{\min}(\Pi_0)}{16\sigma_0^2\gamma_2^2}\right)\right)
\right\}.
\end{align*}
\end{lemma}

\begin{proof}
Fix $\Pi\in\mP_{n,k}\setminus\{\Pi_0\}$ and write
\[
w_\Pi:=\left(I_n-M_{\Pi\bm X}\right)\Pi_0\bm X\bm\beta_0.
\]
Using the identity
\[
M_{(\Pi\bm X,\bm U^*)}
=
M_{\Pi\bm X}+M_{(I_n-M_{\Pi\bm X})\bm U^*},
\]
we have on $\mathcal E(\gamma_1,\gamma_2)$
\begin{align*}
\left\|\left(I_n-M_{(\Pi\bm X,\bm U^*)}\right)\Pi_0\bm X\bm\beta_0\right\|_2^2
&=
\left(1-\phi\left((I_n-M_{\Pi\bm X})\bm U^*,w_\Pi\right)\right)\|w_\Pi\|_2^2\\
&\ge (1-\gamma_1^2)\|w_\Pi\|_2^2.
\end{align*}

Also, on $\{\phi(\bm U^*,\bm U)>1-\gamma_2^2\}$, we have
\[
\left\|\left(I_n-M_{\bm U^*}\right)\bm U\right\|_2\le \gamma_2\|\bm U\|_2,
\]
and since $\range\left(I_n-M_{(\Pi\bm X,\bm U^*)}\right)\subseteq \range\left(I_n-M_{\bm U^*}\right)$ we obtain
\[
\left\|\left(I_n-M_{(\Pi\bm X,\bm U^*)}\right)\bm U\right\|_2
\le
\left\|\left(I_n-M_{\bm U^*}\right)\bm U\right\|_2
\le \gamma_2\|\bm U\|_2.
\]

A direct expansion gives
\begin{align*}
F(\Pi,\bm U^*)-F(\Pi_0,\bm U^*)
&=
\left\|\left(I_n-M_{(\Pi\bm X,\bm U^*)}\right)\Pi_0\bm X\bm\beta_0\right\|_2^2
+2\sigma_0\bm U^\top \left(I_n-M_{(\Pi\bm X,\bm U^*)}\right)\Pi_0\bm X\bm\beta_0\\
&\quad
-\sigma_0^2\bm U^\top\left(M_{(\Pi\bm X,\bm U^*)}-M_{(\Pi_0\bm X,\bm U^*)}\right)\bm U.
\end{align*}
Combining the two bounds above, on $\mathcal E(\gamma_1,\gamma_2)$, the event $\{F(\Pi,\bm U^*)<F(\Pi_0,\bm U^*)\}$ implies
\begin{align}
(1-\gamma_1^2)\|w_\Pi\|_2^2
-2\sigma_0\gamma_2\|\bm U\|_2\|w_\Pi\|_2
-\sigma_0^2\gamma_2^2\|\bm U\|_2^2<0.
\end{align}
If $\|\bm U\|_2\le \frac{1-\gamma_1^2}{4\sigma_0\gamma_2}\|w_\Pi\|_2$, then using $1-\gamma_1^2\le 1$ the left-hand side is at least
$\frac{7}{16}(1-\gamma_1^2)\|w_\Pi\|_2^2>0$, which is a contradiction. Therefore, we have
\[
\|\bm U\|_2^2>\frac{(1-\gamma_1^2)^2}{16\sigma_0^2\gamma_2^2}\|w_\Pi\|_2^2,
\]
and hence either
\[
\|\bm U\|_2^2>\frac{1-\gamma_1^2}{2\sigma_0^2\gamma_2^2}\|w_\Pi\|_2^2\qquad\text{or}\qquad
\|\bm U\|_2^2>\frac{(1-\gamma_1^2)^2}{16\sigma_0^2\gamma_2^2}\|w_\Pi\|_2^2
\]
holds. By definition of $C_{\min}(\Pi_0)$, we have $\|w_\Pi\|_2^2\ge n C_{\min}(\Pi_0)$. Hence, it holds that
\begin{align*}
\bP\left(
F(\Pi,\bm U^*)<F(\Pi_0,\bm U^*)
\mid 
\mathcal E(\gamma_1,\gamma_2)
\right)
&\le
\bP\left(
\|\bm U\|_2^2>\frac{n(1-\gamma_1^2)C_{\min}(\Pi_0)}{2\sigma_0^2\gamma_2^2}
\right)\\
&\quad+
\bP\left(
\|\bm U\|_2^2>\frac{n(1-\gamma_1^2)^2C_{\min}(\Pi_0)}{16\sigma_0^2\gamma_2^2}
\right).
\end{align*}
Since $\|\bm U\|_2^2\sim\chi_n^2$ and $\mathcal E(\gamma_1,\gamma_2)$ depends only on directions,
conditioning on $\mathcal E(\gamma_1,\gamma_2)$ does not change the law of $\|\bm U\|_2^2$.
Applying the $\chi^2$ upper-tail bound (cf. \cite{ghosh2021exponential}) and the union bound over
$\Pi\in\mP_{n,k}\setminus\{\Pi_0\}$, we have the desired result.
\end{proof}


\begin{lemma}\label{lem:max-sim}
Fix $\gamma_1\in(0,1)$. Then, it holds that
\[
\bP\left(
\max_{\Pi\in\mP_{n,k}}
\phi\left((I_n-M_{\Pi\bm X})\bm U^*, (I_n-M_{\Pi\bm X})\Pi_0\bm X\bm\beta_0\right)
\ge \gamma_1^2
\right)
\le
\frac{2}{\pi}|\mP_{n,k}|\left(\arccos(\gamma_1)\right)^{n-p-1}.
\]
\end{lemma}

\begin{proof}
Fix $\Pi\in\mP_{n,k}$ and set $P_\Pi:=I_n-M_{\Pi\bm X}$.
Let $r:=\rank(P_\Pi)=n-p$. Define
\[
Z:=P_\Pi\bm U^*,
\qquad
w:=P_\Pi\Pi_0\bm X\bm\beta_0 \in \range(P_\Pi).
\]
If $w=\bm 0$, then by convention $\phi(Z,w)=0$ (and, in any case, the event
$\{\phi(Z,w)\ge \gamma_1^2\}$ is empty for $\gamma_1>0$), so the desired bound holds trivially.
Hence, assume that $w\neq \bm 0$ and set $e:=w/\|w\|_2$, which is a unit vector in $\range(P_\Pi)$.


Since $P_\Pi$ is an orthogonal projector of rank $r$, there exists an orthogonal matrix
$Q\in\mathbb R^{n\times n}$ such that
\[
Q^\top P_\Pi Q=\begin{pmatrix} I_r & 0\\ 0 & 0\end{pmatrix}.
\]
Write $\bm U^*=Q\begin{pmatrix}G\\H\end{pmatrix}$, where $(G,H)\sim N(0,I_n)$ with
$G\sim N(0,I_r)$ and $H\sim N(0,I_{n-r})$ independent. Then, we have
\[
Z=P_\Pi\bm U^*=Q\begin{pmatrix}G\\0\end{pmatrix}.
\]
In particular, conditional on $\{G\neq 0\}$, the direction $G/\|G\|_2$ is uniform on the unit
sphere $\mathbb S^{r-1}\subset\mathbb R^r$ and is independent of $\|G\|_2$.
Because $Q$ is an isometry and $\range(P_\Pi)=Q(\mathbb R^r\times\{0\})$, it follows that $V:=\frac{Z}{\|Z\|_2}
$ is uniform on the unit sphere in $\range(P_\Pi)$.

Observe that
\[
\phi(Z,w)=\frac{(Z^\top w)^2}{\|Z\|_2^2\|w\|_2^2}
=\left(\frac{Z^\top e}{\|Z\|_2}\right)^2
=(V^\top e)^2.
\]
Hence, we have
\[
\left\{\phi(Z,w)\ge \gamma_1^2\right\}
=
\left\{|V^\top e|\ge \gamma_1\right\}.
\]
Let $\theta:=\arccos(V^\top e)\in[0,\pi]$ be the angle between $V$ and $e$, and set
$\alpha:=\arccos(\gamma_1)\in(0,\pi/2]$. Then $|V^\top e|\ge \gamma_1$ is equivalent to
\[
\theta\in[0,\alpha]\ \cup\ [\pi-\alpha,\pi],
\]
which is the union of two spherical caps of half-angle $\alpha$ centered at $\pm e$.

Because $V$ is uniform on $\mathbb S^{r-1}$, the angle $\theta$ has density proportional to
$\sin^{r-2}(\theta)$, i.e.
\[
\bP(\theta\in dt)=c_r(\sin t)^{r-2}dt,
\qquad
c_r^{-1}=\int_0^\pi (\sin t)^{r-2}dt.
\]
Therefore, with $\alpha:=\arccos(\gamma_1)\in(0,\pi/2]$, we obtain
\[
\bP\left(|V^\top e|\ge \gamma_1\right)
=2\frac{\int_0^\alpha (\sin t)^{r-2}dt}{\int_0^\pi (\sin t)^{r-2}dt}.
\]
Using $\sin t\le t$ for $t\ge 0$ and $\sin t\ge \frac{2}{\pi}t$ for $t\in[0,\pi/2]$, we have
\[
\bP\left(|V^\top e|\ge \gamma_1\right)
\le
2\frac{\int_0^\alpha t^{r-2}dt}{2\int_0^{\pi/2} \left(\frac{2}{\pi}t\right)^{r-2}dt}
=
\frac{2}{\pi}\alpha^{r-1}
=
\frac{2}{\pi}\left(\arccos(\gamma_1)\right)^{r-1}.
\]


Since $\{\phi(Z,w)\ge \gamma_1^2\}=\{|V^\top e|\ge \gamma_1\}$, this proves that
\[
\bP\left(\phi\left(P_\Pi\bm U^*,P_\Pi\Pi_0\bm X\bm\beta_0\right)\ge \gamma_1^2\right)
\le \frac{2}{\pi}\left(\arccos(\gamma_1)\right)^{r-1}
\]
for fixed $\Pi$.
Finally, applying a union bound over $\Pi\in\mP_{n,k}$ completes the proof.
\end{proof}



\begin{lemma}\label{lem:inconsistency}
Let $\gamma\in(0,1/4]$ and let $\widehat{\Pi}$ be the optimizer in \eqref{eq:repro_Pi} based on a single repro sample $\bm U^*$. Then, we have
\[
\bP_{(\bm U^*,\bm U)}\left(
\widehat{\Pi}\neq \Pi_0 \mid \phi(\bm U^*,\bm U)>1-\gamma^2
\right)
\le \Delta(\gamma).
\]
\end{lemma}

\begin{proof}
Fix any $\gamma\in(0,1/4]$ and set $\gamma_2:=\gamma$.  Define 
\[
\gamma_1
:=
\sqrt{1-\gamma\log\frac{e}{\gamma}}
\in(0,1),
\]
where the inclusion $\gamma_1\in(0,1)$ holds since the function
$f(\gamma):=\gamma\log(e/\gamma)$ is increasing on $(0,1)$ and hence
$f(\gamma)\le f(1/4)=\frac14\log(4e)<1$.

Let us define the event
\[
\mathcal{E}
:=
\left\{
\max_{\Pi\in\mP_{n,k}}
\phi\left(
\bm P_\Pi \bm U^*,
\bm P_\Pi \Pi_0 \bm X\beta_0
\right)
<
\gamma_1^2
\right\}.
\]
By Lemma~\ref{lem:proj-diff} together with Lemma~\ref{lem:F-union}, on the event
$\mathcal{E}\cap\{\phi(\bm U^*,\bm U)>1-\gamma_2^2\}$ we have $\widehat{\Pi}=\Pi_0$.
Therefore, it holds that
\begin{align}
\bP_{(\bm U^*,\bm U)}\left(
\widehat{\Pi}\neq \Pi_0
\mid 
\phi(\bm U^*,\bm U)>1-\gamma_2^2
\right)
&\le
\bP\left(
\mathcal{E}^c
\mid 
\phi(\bm U^*,\bm U)>1-\gamma_2^2
\right)
\nonumber\\
&\quad+
\bP\left(
\widehat{\Pi}\neq \Pi_0
\mid
\mathcal{E}\cap\{\phi(\bm U^*,\bm U)>1-\gamma_2^2\}
\right).
\label{eq:inconsistency-split}
\end{align}

We first bound the term involving $\mathcal{E}^c$.
The joint law of $(\bm U^*,\bm U)$ is invariant under simultaneous orthogonal
transformations, and the conditioning event
$\{\phi(\bm U^*,\bm U)>1-\gamma_2^2\}$ depends only on the angle between
$\bm U^*$ and $\bm U$. Hence, the conditional law of $\bm U^*$ given
$\{\phi(\bm U^*,\bm U)>1-\gamma_2^2\}$ remains spherically symmetric in
$\R^n$. Since $\phi(\cdot,\cdot)$ is scale-invariant, we may apply
Lemma~\ref{lem:max-sim} under this conditioning to obtain
\[
\bP\left(
\mathcal{E}^c
 \mid 
\phi(\bm U^*,\bm U)>1-\gamma_2^2
\right)
\le
\frac{2}{\pi}|\mP_{n,k}|
\left(\arccos(\gamma_1)\right)^{n-p-1},
\]
Noting that
\[
\arccos(\gamma_1)
=
\arccos\left(\sqrt{1-\gamma\log\frac{e}{\gamma}}\right)
=
\arcsin\left(\sqrt{\gamma\log\frac{e}{\gamma}}\right)
\le
\frac{\pi}{2}\sqrt{\gamma\log\frac{e}{\gamma}},
\]
so we obtain
\[
\bP\left(
\mathcal{E}^c
 \mid 
\phi(\bm U^*,\bm U)>1-\gamma_2^2
\right)
\le
\left(\frac{\pi}{2}\right)^{n-p-2}|\mP_{n,k}|
\left(
\gamma\log\frac{e}{\gamma}
\right)^{\frac{n-p-1}{2}}.
\]

For the second term on the right-hand side of \eqref{eq:inconsistency-split},
Lemma~\ref{lem:F-union} applied with $(\gamma_1,\gamma_2)$ yields
\begin{align*}
&\bP\left(
\widehat{\Pi}\neq \Pi_0
\mid
\mathcal{E}\cap\{\phi(\bm U^*,\bm U)>1-\gamma_2^2\}
\right)\\&
\le
\sum_{m=n-k}^n\binom{n}{m}D_{n-m}
\left\{
\exp\left(
-\frac{n}{2}\Psi\left(
\max\left\{1,
\frac{C_{\min}(\Pi_0)(1-\gamma_1^2)}{2\sigma_0^2\gamma_2^2}
\right\}
\right)
\right)\right.\\&
+
\left.\exp\left(
-\frac{n}{2}\Psi\left(
\max\left\{1,
\frac{C_{\min}(\Pi_0)(1-\gamma_1^2)^2}{16\sigma_0^2\gamma_2^2}
\right\}
\right)
\right)
\right\}.    
\end{align*}

Finally, substituting $1-\gamma_1^2=\gamma\log(e/\gamma)$ and $\gamma_2=\gamma$
and combining the two bounds gives \eqref{eq:Delta-gamma}.
\end{proof}



\begin{lemma}\label{lem:proj-diff}
For any $\Pi\in\mP_n$, define
\[
v_0:=(I_n-M_{\bm X})\bm U^*,
\qquad
v_\Pi:=(I_n-M_{\bm X})\Pi^\top \bm U^* .
\]
Then, it follows that
\begin{equation}
\left\|\Pi M_{(\bm X,\bm U^*)} - M_{(\Pi \bm X,\bm U^*)}\right\|_{\mathrm{op}}
=
\left\|M_{v_0}-M_{v_\Pi}\right\|_{\mathrm{op}} .
\label{eq:PiPminusM}
\end{equation}
In particular, whenever $v_0$ and $v_\Pi$ are nonzero, we have
\begin{equation}
\left\|M_{v_0}-M_{v_\Pi}\right\|_{\mathrm{op}}
\le
\frac{2\|v_0-v_\Pi\|_2}{\min\{\|v_0\|_2,\|v_\Pi\|_2\}}
\le
\frac{2\|\bm U^*-\Pi^\top \bm U^*\|_2}{\min\{\|v_0\|_2,\|v_\Pi\|_2\}} .
\label{eq:rank1-diff}
\end{equation}
\end{lemma}

\begin{proof}
By orthogonal invariance of projectors, we have
\[
M_{(\Pi \bm X,\bm U^*)}=\Pi M_{(\bm X,\Pi^\top \bm U^*)}\Pi^\top.
\]
Hence, it holds that
\[
\Pi^\top\left(\Pi M_{(\bm X,\bm U^*)}-M_{(\Pi \bm X,\bm U^*)}\right)\Pi
= M_{(\bm X,\bm U^*)}-M_{(\bm X,\Pi^\top \bm U^*)}.
\]
For any vector $\bm{z} \in \R^n$, the subspaces $\mathrm{span}(\bm X)$ and $\mathrm{span}\left((I_n-M_{\bm X})z\right)$ are orthogonal, and hence we have
\[
M_{(\bm X,\bm{z})} = M_{\bm X} + M_{(I_n-M_{\bm X})\bm{z}}.
\]
Applying this with $\bm{z}=\bm U^*$ and $z=\Pi^\top \bm U^*$ gives
\[
M_{(\bm X,\bm U^*)}-M_{(\bm X,\Pi^\top \bm U^*)}
= \left\{M_{\bm X}+M_{v_0}\right\}-\left\{M_{\bm X}+M_{v_\Pi}\right\}
= M_{v_0}-M_{v_\Pi}.
\]
Conjugation by the orthogonal matrix $\Pi$ preserves the operator norm, yielding \eqref{eq:PiPminusM}.

For \eqref{eq:rank1-diff}, write $\hat u:=u/\|u\|_2$ for $u\neq 0$. Then, the following holds:
\[
\|M_{v_0}-M_{v_\Pi}\|_{\mathrm{op}}
=\left\|\hat v_0\hat v_0^\top-\hat v_\Pi\hat v_\Pi^\top\right\|_{\mathrm{op}}
=\sqrt{1-(\hat v_0^\top \hat v_\Pi)^2}
\le \|\hat v_0-\hat v_\Pi\|_2
\le \frac{2\|v_0-v_\Pi\|_2}{\min\{\|v_0\|_2,\|v_\Pi\|_2\}}.
\]
Finally, due to $\|I_n-M_{\bm X}\|_{\mathrm{op}}\le 1$, we obtain $
\|v_0-v_\Pi\|_2
=\|(I_n-M_{\bm X})(\bm U^*-\Pi^\top \bm U^*)\|_2
\le \|\bm U^*-\Pi^\top \bm U^*\|_2$
which gives the last inequality in \eqref{eq:rank1-diff}.
\end{proof}

\begin{lemma}\label{lem:perm-diff-finite}
Fix a set $\mS \subset \mP_{n,k}$ with cardinality $|\mS|=m$. Let $\bm U^*\sim\Normalp{0}{I_n}$. It holds that, with probability at least $1-\xi$,
\begin{equation}\label{eq:sup-diff}
\max_{\Pi\in\mS} \left\|\bm U^*-\Pi^\top\bm U^*\right\|_2
\le
2\left(k+2\sqrt{k\log\frac{m}{\xi}}+2\log\frac{m}{\xi}\right)^{1/2}.
\end{equation}
\end{lemma}


\begin{proof}
Fix $\Pi\in\mP_{n,k}$ and set $S=\{i:\Pi(i)\neq i\}$, so that $|S|\le k$. Then, we have
\[
\|\bm{U}^*-\Pi^\top \bm{U}^*\|_2^2
=\sum_{i\in S}({U}^*_i-{U}^*_{\Pi(i)})^2
\le 2\sum_{i\in S}({U}^*_i)^2+2\sum_{i\in S}({U}^*_{\Pi(i)})^2
\le 4\sum_{i\in S}({U}^*_i)^2.
\]
For any fixed $T\subset[n]$ with $|T|=k$, the sum $\sum_{i\in T}({U}^*_i)^2$ has the $\chi^2_k$ distribution. The upper‑tail bound in Lemma 1 of \cite{laurent2000adaptive} gives that, for any $t>0$,
\[
\bP \left\{\sum_{i\in T}({U}^*_i)^2  \ge k + 2\sqrt{k t} + 2t\right\}  \le e^{-t}.
\]
Taking a union bound over the $m$ choices $T$ and choosing $t=\log \left(m/\xi\right)$ yields that
\[
\bP \left\{
\max_{\Pi \in \mS}\sum_{i\in T}({U}^*_i)^2
 \le
k + 2\sqrt{k \left( \log m+\log(1/\xi) \right)} + 2\left( \log m +\log(1/\xi) \right)
\right\}
 \ge 1-\xi.
\] Taking square roots gives \eqref{eq:sup-diff}.
\end{proof}

\begin{lemma}\label{lem:denom-lb}
Fix a set $\mS \subset \mP_{n,k}$ with cardinality $|\mS|=m$. Assume $\rank(\bm X)=p$ with $n-p\ge 1$. Define
\[
v_0:=(I_n-M_{\bm X}) \bm U^*,
\qquad
v_\Pi:=(I_n-M_{\bm X}) \Pi^\top \bm U^*,
\]
where $\bm U^*\sim\Normalp{0}{I_n}$. Then, for any $\xi\in(0,1)$, with probability at least $1-\xi$, we have
\begin{equation}
\min_{\Pi\in\mS} \min\{\|v_0\|_2,\|v_\Pi\|_2\}
\ge
\sqrt{(n-p)-2\sqrt{(n-p)\log(1/\xi)}}
-\sup_{\Pi\in\mS } \|\bm U^*-\Pi^\top \bm U^*\|_2 .
\label{eq:denom-lb}
\end{equation}
\end{lemma}

\begin{proof}
Since $M_{\bm X}$ projects onto $\mathrm{span}(\bm X)$ with rank $p$, we have
$\|v_0\|_2^2=\|(I_n-M_{\bm X})\bm U^*\|_2^2\sim\chi^2_{n-p}$. Lemma 1 in \cite{laurent2000adaptive} implies that, with probability at least $1-\xi$, we have
\[
\|v_0\|_2  \ge  \sqrt{ (n-p)-2\sqrt{(n-p)\log(1/\xi)}}.
\]
For any $\Pi\in\mS \subset \mP_{n,k}$, the inequalities hold:
\[
\|v_\Pi\|_2
=\|(I_n-M_{\bm X})\Pi^\top \bm U^*\|_2
 \ge \|v_0\|_2 - \|(I_n-M_{\bm X})(\Pi^\top \bm U^*-\bm U^*)\|_2
 \ge \|v_0\|_2 - \|\bm U^*-\Pi^\top \bm U^*\|_2,
\]
using $\|I_n-M_{\bm X}\|_{\mathrm{op}}\le 1$ in the last step. Hence, for each $\Pi$, we have
\[
\min\{\|v_0\|_2,\|v_\Pi\|_2\} \ge \|v_0\|_2 - \|\bm U^*-\Pi^\top \bm U^*\|_2.
\]
Taking the infimum over $\Pi\in\mS \subset \mP_n$ on the left and the supremum on the right gives \eqref{eq:denom-lb}.
\end{proof}

\begin{proposition}\label{prop:finite-path-op}
Under the event in Lemma \ref{lem:perm-diff-finite} and Lemma \ref{lem:denom-lb}, whenever it holds that
\begin{equation}\label{eq:positivity-path}
\sqrt{ (n-p)-2\sqrt{(n-p)\log(1/\xi)}} >2\sqrt{k+2\sqrt{k\log\frac{m}{\xi}  }+2\log\frac{m}{\xi}},
\end{equation}
we have
\begin{align*}\label{eq:op-bound-finite}
& \max_{\Pi\in\mS} \left\|\Pi M_{(\bm X,\bm U^*)}-M_{(\Pi\bm X,\bm U^*)}\right\|_{\mathrm{op}}
 \\ & \le
\frac{4\sqrt{k+2\sqrt{k\log\frac{m}{\xi}}+2\log\frac{m}{\xi}}}{\sqrt{ (n-p)-2\sqrt{(n-p)\log(1/\xi)}}-2\sqrt{k+2\sqrt{k\log\frac{m}{\xi}}+2\log\frac{m}{\xi}}}.
\end{align*}
\end{proposition}

\begin{proof}
By Lemma \ref{lem:proj-diff} it holds thatt
$
\|\Pi M_{(\bm X,\bm U^*)}-M_{(\Pi\bm X,\bm U^*)}\|_{\mathrm{op}}
\le 2\|v_0-v_\Pi\|_2/\min\{\|v_0\|_2,\|v_\Pi\|_2\}.
$
Then, applying Lemma \ref{lem:perm-diff-finite} and \ref{lem:denom-lb} yields the conclusion.

\end{proof}


\begin{lemma}\label{lem:poly-tail-finite}
Let $Z\sim\Normal(0,1)$, $\mu\in\mathbb R$ and $\sigma\ge 0$. For any $\xi\in(0,1)$,
with probability at least $1-\xi$, we obtain
\[
\big|(\mu+\sigma Z)^2-(\mu^2+\sigma^2)\big|
\le 2\sigma|\mu|\sqrt{2\log(4/\xi)}+2\sigma^2\left(\sqrt{\log(4/\xi)}+\log(4/\xi)\right).
\]
\end{lemma}

\begin{proof}
Write $(\mu+\sigma Z)^2-(\mu^2+\sigma^2)=2\mu\sigma Z+\sigma^2(Z^2-1)$. By the Gaussian tail bound,
$|Z|\le \sqrt{2\log(4/\xi)}$ holds with probability at least $1-\xi/2$.
Moreover, since $Z^2\sim\chi^2_1$, Lemma 1 in  \cite{laurent2000adaptive} yields
\[
\bP\left(|Z^2-1|>2(\sqrt{t}+t)\right)\le 2e^{-t},\qquad t>0.
\]
Taking $t=\log(4/\xi)$ gives $|Z^2-1|\le 2(\sqrt{\log(4/\xi)}+\log(4/\xi))$ with probability at least $1-\xi/2$.
The claim follows by a union bound.
\end{proof}

\begin{lemma}\label{lem:Y-and-diag}
Let $\bm U^*\sim \Normal(0,I_n)$ be independent of $\bm U\sim \Normal(0,I_n)$ and define
\[
\bm Y=\Pi_0\bm X\bm\beta_0+\sigma_0\bm U,
\qquad 
\bm m:=M_{(\bm X,\bm U^*)}\bm Y .
\]
Write $m_i$ for the $i$-th coordinate of $\bm m$.
Then, with probability at least $1-\xi$, it holds that
\begin{equation}\label{eq:Y-upper-finite}
\|\bm Y\|_2^2\ \le\
\|\Pi_0\bm X\bm\beta_0\|_2^2
+2\sigma_0\|\Pi_0\bm X\bm\beta_0\|_2\sqrt{2\log\frac{2}{\xi}}
+\sigma_0^2\left(n+2\sqrt{n\log\frac{2}{\xi}}+2\log\frac{2}{\xi}\right).
\end{equation}
Also, with probability at least $1-\xi$, it holds that
\begin{align}\label{eq:diag-sq-finite}
&\max_{1\le i\le n}\left|(Y_i-m_i)^2-\bE\left[(Y_i-m_i)^2\mid \bm U^*\right]\right|\\ &\le 
2\sigma_0\left\|(I_n-M_{(\bm X,\bm U^*)})\Pi_0\bm X\bm\beta_0\right\|_\infty\sqrt{2\log\frac{8n}{\xi}}
+2\sigma_0^2\left(\sqrt{\log\frac{8n}{\xi}}+\log\frac{8n}{\xi}\right).
\end{align}
\end{lemma}

\begin{proof}
For \eqref{eq:Y-upper-finite}, write $\bm Y=\Pi_0\bm X\bm\beta_0+\sigma_0\bm U$ and expand
\[
\|\bm Y\|_2^2
=
\|\Pi_0\bm X\bm\beta_0\|_2^2
+2\sigma_0\langle \Pi_0\bm X\bm\beta_0,\bm U\rangle
+\sigma_0^2\|\bm U\|_2^2.
\]
Since $\langle \Pi_0\bm X\bm\beta_0,\bm U\rangle\sim \Normal\left(0,\|\Pi_0\bm X\bm\beta_0\|_2^2\right)$ and $\|\bm U\|_2^2\sim \chi_n^2$, standard Gaussian and chi-square tail bounds yield \eqref{eq:Y-upper-finite}.

For \eqref{eq:diag-sq-finite}, condition on $\bm U^*$ and set
\[
\bm r:=\bm Y-\bm m=(I_n-M_{(\bm X,\bm U^*)})\bm Y,
\qquad 
a_i:=(I_n-M_{(\bm X,\bm U^*)})e_i.
\]
It gives $r_i=Y_i-m_i=a_i^\top\bm Y=\mu_i+\sigma_0 a_i^\top \bm U$ where
$\mu_i:=a_i^\top \Pi_0\bm X\bm\beta_0$.
Define $s_i:=\|a_i\|_2$ so that $a_i^\top \bm U\sim \Normal(0,s_i^2)$ and $s_i\le 1$.
Applying Lemma \ref{lem:poly-tail-finite} to the sub-exponential variables $(r_i)^2$ and taking a union bound over $i\in[n]$ gives that, with probability at least $1-\xi$,
\[
\max_{1\le i\le n}\left|(Y_i-m_i)^2-\bE\left[(Y_i-m_i)^2\mid \bm U^*\right]\right|
\le
2\sigma_0\max_{1\le i\le n}|\mu_i|\sqrt{2\log\frac{8n}{\xi}}
+2\sigma_0^2\left(\sqrt{\log\frac{8n}{\xi}}+\log\frac{8n}{\xi}\right).
\]
This yields \eqref{eq:diag-sq-finite}.
\end{proof}




\section{Proofs of main results}
\subsection{Proof of Proposition \ref{prop:identification}}\label{proof:prop identification}
\begin{proof}
Since $\bm Y^{\obs}=\Pi_0\bm X\bm\beta_0+\sigma_0\bm u^{\rel}$, we have
\[
\left\|(I_n-M_{(\Pi_0\bm X,\bm u^{\rel})})\bm Y^{\obs}\right\|_2^2=0,
\]
so the minimum value in \eqref{eq:repro_ideal_Pi} is $0$.

Let $\tilde\Pi\in\mP_{n,k}$ satisfy
$\|(I_n-M_{(\tilde\Pi\bm X,\bm u^{\rel})})\bm Y^{\obs}\|_2^2=0$.
Then, it holds that $\bm Y^{\obs}\in\range([\tilde\Pi\bm X,\bm u^{\rel}])$ so there exist
$(\tilde{\bm\beta},\tilde\sigma)$ such that
\[
\bm Y^{\obs}=\tilde\Pi\bm X\tilde{\bm\beta}+\tilde\sigma\bm u^{\rel}.
\]
Subtracting the true representation gives
\begin{equation}\label{eq:oracle-basic-eq}
\Pi_0\bm X\bm\beta_0-\tilde\Pi\bm X\tilde{\bm\beta}=(\tilde\sigma-\sigma_0)\bm u^{\rel}.
\end{equation}

Define $P:=\Pi_0^\top\tilde\Pi$ and $\tilde{\bm u}:=\Pi_0^\top\bm u^{\rel}$.
Multiplying \eqref{eq:oracle-basic-eq} by $\Pi_0^\top$ yields
\begin{equation}\label{eq:oracle-permuted}
\bm X\bm\beta_0- P\bm X\tilde{\bm\beta}=(\tilde\sigma-\sigma_0)\tilde{\bm u}.
\end{equation}
Note that $P\in\mP_{n,2k}$ , hence its fixed-point set
$T:=\{i\in[n]:P(i)=i\}$ satisfies $|T|\ge n-2k\ge p$.

Restrict \eqref{eq:oracle-permuted} to the coordinates in $T$. On the set $T$, we have
$(P\bm X\tilde{\bm\beta})_T=\bm X_T\tilde{\bm\beta}$, so it holds that
\begin{equation}\label{eq:oracle-on-T}
\bm X_T(\bm\beta_0-\tilde{\bm\beta})=(\tilde\sigma-\sigma_0)\tilde{\bm u}_T.
\end{equation}

We now show that $\tilde\sigma=\sigma_0$ and $\tilde{\bm\beta}=\bm\beta_0$ almost surely (conditional on $\bm X$).
Fix any non-identity $P\in\mP_{n,2k}$ with fixed-point set $T$.
Under Assumption~\ref{ass:generic-design}, $\rank(\bm X_T)=p$ almost surely.

If $|T|>p$, then $\range(\bm X_T)$ is a proper $p$-dimensional subspace of $\R^{|T|}$.
Since $\tilde{\bm u}_T\sim\Normal(0,I_{|T|})$ has a Lebesgue density, we have
$\bP(\tilde{\bm u}_T\in\range(\bm X_T))=0$.
Therefore \eqref{eq:oracle-on-T} cannot hold with $\tilde\sigma\ne\sigma_0$; hence $\tilde\sigma=\sigma_0$ and then
$\tilde{\bm\beta}=\bm\beta_0$.

If $|T|=p$, then $\bm X_T$ is square and invertible almost surely, and \eqref{eq:oracle-on-T} gives
\[
\tilde{\bm\beta}=\bm\beta_0-(\tilde\sigma-\sigma_0)\bm X_T^{-1}\tilde{\bm u}_T.
\]
Plugging this into \eqref{eq:oracle-permuted} on any index $i\notin T$ yields an equation of the form
\[
(\tilde\sigma-\sigma_0)\left\{\tilde u_i-a_i^\top \tilde{\bm u}_T\right\}=b_i,
\]
where $a_i$ and $b_i$ are nonrandom given $\bm X$ and $(P,\bm\beta_0)$.
Since $P\neq I_n$ implies $|[n]\setminus T|\ge 2$, using two such indices shows that the event
$\tilde\sigma\ne\sigma_0$ forces a nontrivial algebraic relation among independent Gaussian coordinates
$\tilde u_i,\tilde u_j,\tilde{\bm u}_T$, which has probability zero. Hence $\tilde\sigma=\sigma_0$ and then
$\tilde{\bm\beta}=\bm\beta_0$ also in this boundary case.

Taking a union bound over the finite set of possible $P\in\mP_{n,2k}$ yields that,
\[
\tilde\sigma=\sigma_0
\quad\text{and}\quad
P\bm X\bm\beta_0=\bm X\bm\beta_0
\]
almost surely. Finally, since the entries of $\bm X\bm\beta_0$ are pairwise distinct by assumption,
the only permutation fixing $\bm X\bm\beta_0$ is the identity, so we have $P=I_n$ and $\tilde\Pi=\Pi_0$.
This proves uniqueness of $\Pi_0$ in \eqref{eq:repro_ideal_Pi}.
\end{proof}

\subsection{Proof of Theorem \ref{thm:recovery}}

\begin{proof}
For a nonnegative integer $m$, let $D_m$ denote the number of derangements of $m$ elements, i.e.,
the number of permutations of $\{1,\dots,m\}$ with no fixed points (see Lemma~\ref{lem:k-sparse permutation}).

Define $\Psi(x):=x-1-\log x$ for $x>0$. Fix any $\gamma\in(0,1/4]$ and recall that
\[
\mC^{(L)}=\{\hat\Pi_\ell:\ell=1,\dots,L\}.
\]
Define
\begin{align}
\label{eq:Delta-gamma}
\Delta(\gamma)
:=
&\sum_{m=n-k}^n\binom{n}{m}D_{n-m}
\Bigg\{
\exp\Bigg(
-\frac{n}{2}\Psi\left(
\max\left\{1,
\frac{C_{\min}(\Pi_0)\log(e/\gamma)}{2\sigma_0^2\gamma}
\right\}
\right)
\Bigg)
\nonumber\\
&\qquad\qquad+
\exp\Bigg(
-\frac{n}{2}\Psi\left(
\max\left\{1,
\frac{C_{\min}(\Pi_0)\log^2(e/\gamma)}{16\sigma_0^2}
\right\}
\right)
\Bigg)
\Bigg\}
\nonumber\\
&\quad+
\left(\frac{\pi}{2}\right)^{n-p-2}\lvert\mP_{n,k}\rvert
\left(\gamma\log\frac{e}{\gamma}\right)^{\frac{n-p-1}{2}}.
\end{align}

Also, define the alignment events as
\[
A_\ell := \{\phi(\bm U_\ell^*,\bm U)>1-\gamma^2\},
\qquad \ell=1,\dots,L,
\]
and the (censored) stopping index
\[
\tau := \min\{\ell\in\{1,\dots,L\}: A_\ell\},
\qquad \text{with the convention }\tau=\infty \text{ if }A_1,\dots,A_L\text{ all fail.}
\]
If $\Pi_0\notin\mC^{(L)}$, then either $\tau=\infty$ or $\hat\Pi_\tau\neq\Pi_0$ holds when $\tau\le L$. Hence, we have
\begin{equation}\label{eq:recovery-split}
\bP(\Pi_0\notin\mC^{(L)})
\le \bP(\tau=\infty)+\bP(\hat\Pi_\tau\neq\Pi_0,\ \tau\le L).
\end{equation}

Lemma~\ref{lem:similarity measure} (applied with $\gamma_2=\gamma$) gives
\[
\bP(\tau=\infty)
=\bP\left(\bigcap_{\ell=1}^L A_\ell^c\right)
\le \left(1-\frac{\gamma^{n-1}}{n-1}\right)^L.
\]

Define $B_\ell := \{\hat\Pi_\ell\neq \Pi_0\}$ for $\ell=1,\dots,L$.
Conditionally on $\bm U$, it holds that $\bm U_1^*,\dots,\bm U_L^*$ are i.i.d.\ and, for each $\ell$,
both $A_\ell$ and $B_\ell$ are measurable functions of $(\bm U,\bm U_\ell^*)$.
Hence the pairs $(A_\ell,B_\ell)$ are i.i.d.\ across $\ell$ given $\bm U$.
Define
\[
q(\bm U):=\bP(A_1\mid \bm U),\qquad r(\bm U):=\bP(A_1\cap B_1\mid \bm U).
\]
By the definition of $\tau$, for $t=1,\dots,L$, we have
\[
\{\tau=t\}=\left(\bigcap_{j=1}^{t-1}A_j^c\right)\cap A_t,
\]
so by conditional independence, we have
\begin{align*}
\bP(B_\tau,\ \tau\le L\mid \bm U)
&=\sum_{t=1}^L \bP(\tau=t,\ B_t\mid \bm U)\\
&=\sum_{t=1}^L \bP\left(\bigcap_{j=1}^{t-1}A_j^c,\ A_t\cap B_t\mid \bm U\right)\\
&=\sum_{t=1}^L (1-q(\bm U))^{t-1} r(\bm U)\\
&= r(\bm U)\frac{1-(1-q(\bm U))^L}{q(\bm U)}
= \bigl(1-(1-q(\bm U))^L\bigr)\bP(B_1\mid \bm U, A_1)\\
&\le \bP(B_1\mid \bm U, A_1).
\end{align*}

Note that $q(\bm U)=\bP(A_1\mid \bm U)$ is a constant (does not depend on $\bm U$) by rotational
invariance of $\bm U_1^*\sim N(0,I_n)$ and the fact that $A_1$ depends on $\bm U_1^*,\bm U$
only through their angle. This yields that
\[
\bP(B_\tau,\tau\le L)
=\sum_{t=1}^L (1-q)^{t-1}\bP(A_1\cap B_1)
=(1-(1-q)^L)\bP(B_1\mid A_1)
\le \bP(B_1\mid A_1)\le \Delta(\gamma),
\]
where the last inequality is from Lemma~\ref{lem:inconsistency}.

Therefore, we have
\[
\bP(\hat\Pi_\tau\neq\Pi_0,\ \tau\le L)\le \Delta(\gamma).
\]

Combining the bounds with \eqref{eq:recovery-split} gives that, for any $\gamma\in(0,1/4]$,
\[
\bP(\Pi_0\notin\mC^{(L)})
\le \Delta(\gamma)+\left(1-\frac{\gamma^{n-1}}{n-1}\right)^L.
\]

Now set 
\begin{equation}\label{eq:gamma-L}
  \gamma_L := \min\left\{\frac14, \left(\frac{(n-1)\log(e L)}{L}\right)^{1/(n-1)}\right\}.  
\end{equation}

Since $(1-x)^L\le e^{-Lx}$ for $x\in[0,1]$, for $L$ large enough, we have
\[
\left(1-\frac{\gamma_L^{n-1}}{n-1}\right)^L
\le \exp\left(-\frac{L\gamma_L^{n-1}}{n-1}\right)
\le \exp(-\log(eL))=\frac{1}{eL}.
\]
Moreover, by inspection of \eqref{eq:Delta-gamma}, the only term in $\Delta(\gamma_L)$ contributing at polynomial order in $L$
is $\left(\gamma_L\log(e/\gamma_L)\right)^{(n-p-1)/2}$, and hence we obtain
\[
\Delta(\gamma_L)
=
O\left(\left(\gamma_L\log\frac{e}{\gamma_L}\right)^{\frac{n-p-1}{2}}\right)
=
O\left(\frac{(\log L)^{\frac{n-p}{2}}}{L^{\frac{n-p-1}{2(n-1)}}}\right).
\]
This yields the claim.
\end{proof}



\subsection{Proof of Theorem \ref{thm:p-value}}
We prove Proposition~\ref{thm:p-value_without_delta} first, and then conclude Theorem~\ref{thm:p-value} by adding the localization error.

\begin{proposition}[Type-I error control for the sparsity test]\label{thm:p-value_without_delta}
For any $k_0\le k$ and $\theta_0=(\Pi_0,\bm\beta_0,\sigma_0)$ with $\Pi_0\in\mP_{n,k_0}$, we have
\begin{equation}\label{eq:conditional-validity-k0}
\bP_{\left(\mathcal{U}^L, \bm{U}\right)}\left(
D(\bm Y_{\theta_0})\notin \hat{B}^{(k_0)}_{1-\alpha}\ \text{ and }\ \Pi_0\in \mC^{(L)}
\mid 
S_{1,\Pi_0}(\bm Y_{\theta_0}),S_{2,\Pi_0}(\bm Y_{\theta_0})
\right)\le \alpha.
\end{equation}
\end{proposition}

\begin{proof}
Fix $\theta_0=(\Pi_0,\bm\beta_0,\sigma_0)$ with $\Pi_0\in\mP_{n,k_0}$ and write
\[
s_{1,0}:=S_{1,\Pi_0}(\bm Y_{\theta_0}),
\qquad
s_{2,0}:=S_{2,\Pi_0}(\bm Y_{\theta_0}).
\]
By the conditional representation \eqref{eq:conditional-representation} (with $\Pi=\Pi_0$), conditionally on $(s_{1,0},s_{2,0})$ we have
\[
\bm Y_{\theta_0}\mid (s_{1,0},s_{2,0})
\ \overset{d}{=}\
s_{1,0}
+
s_{2,0}\frac{(I_n-M_{\Pi_0\bm X})\widetilde{\bm U}}{\|(I_n-M_{\Pi_0\bm X})\widetilde{\bm U}\|_2},
\]
where $\widetilde{\bm U}\sim\Normal(0,I_n)$ is independent of $(s_{1,0},s_{2,0})$.

On the event $\{\Pi_0\in \mC^{(L)}\}$, by definition of the localized null set (see Section~\ref{subsubsection:testing permutation}) we have
$\Pi_0\in\hat{\mP}_{n,k_0}$, and hence the composite critical value in \eqref{eq:critical-value-composite} satisfies
\[
\widehat c_{1-\alpha}^{(k_0)}
=\max_{\Pi\in \hat{\mP}_{n,k_0}}\widehat c_{1-\alpha}(\Pi)
\ \ge\
\widehat c_{1-\alpha}(\Pi_0),
\]
where $\widehat c_{1-\alpha}(\Pi_0)$ is the conditional $(1-\alpha)$-quantile of $D(\bm Y)$ given
$\left(S_{1,\Pi_0}(\bm Y),S_{2,\Pi_0}(\bm Y)\right)=(s_{1,0},s_{2,0})$
(as defined in Section~\ref{subsubsection:testing permutation}).

Therefore, we have
\begin{align*}
&\bP_{\left(\mathcal{U}^L, \bm{U}\right)}\left(
D(\bm Y_{\theta_0})\notin \hat B^{(k_0)}_{1-\alpha}\ \text{ and }\ \Pi_0\in \mC^{(L)}
\mid s_{1,0},s_{2,0}
\right)\\
&\qquad\le
\bP_{\left(\mathcal{U}^L, \bm{U}\right)}\left(
D(\bm Y_{\theta_0})>
\widehat c_{1-\alpha}(\Pi_0)
\mid s_{1,0},s_{2,0}
\right)
 \le \alpha,
\end{align*}
where the last inequality follows from the defining property of the conditional $(1-\alpha)$-quantile.
\end{proof}

\begin{proof}[Proof of Theorem~\ref{thm:p-value}]
Let $\theta_0=(\Pi_0,\bm\beta_0,\sigma_0)$ with $\Pi_0\in\mP_{n,k_0}$.
By a union bound, we have
\begin{align*}
\bP_{\left(\mathcal{U}^L,\bm U\right)}\left(D(\bm Y_{\theta_0})\notin \hat{B}^{(k_0)}_{1-\alpha}\right)
&\le
\bP_{\left(\mathcal{U}^L,\bm U\right)}\left(D(\bm Y_{\theta_0})\notin \hat{B}^{(k_0)}_{1-\alpha}\ \text{and}\ \Pi_0\in\mC^{(L)}\right)
+
\bP_{\left(\mathcal{U}^L,\bm U\right)}\left(\Pi_0\notin\mC^{(L)}\right)\\
&\le \alpha+\delta_L,
\end{align*}
where the second inequality uses Proposition~\ref{thm:p-value_without_delta} for the first term and the localization guarantee
$\bP(\Pi_0\notin\mC^{(L)})\le \delta_L$ for the second term (see Theorem~\ref{thm:recovery}).
\end{proof}



\subsection{Proof of Theorem \ref{thm:confidence}}
\begin{proof}
By definition, the confidence set is given by
\[
\Gamma_{\alpha}^{\mathrm{coef}}(\bm{Y}^{\obs})
= \bigcup_{\Pi\in \mC^{(L)}}\left\{\bm{\beta}\in \R^p:\ T(\bm{Y}^{\obs},(\Pi,\bm{\beta}))\le F_{p,n-p}^{-1}(\alpha)\right\}.
\]
Hence, on the event $\{\Pi_0\in \mC^{(L)}\}$ we have
\[
\{\bm{\beta}_0\notin \Gamma_{\alpha}^{\mathrm{coef}}(\bm{Y}^{\obs})\}\cap\{\Pi_0\in \mC^{(L)}\}
\subseteq
\left\{T(\bm{Y}^{\obs},(\Pi_0,\bm{\beta}_0))> F_{p,n-p}^{-1}(\alpha)\right\}.
\]
Under $\bm{Y}^{\obs}=\Pi_0\bm X\bm\beta_0+\sigma_0\bm U$ with $\bm U\sim \Normal(0,I_n)$ and $\Pi_0\bm X$ having rank $p$, we have
$T(\bm{Y}^{\obs},(\Pi_0,\bm{\beta}_0))\sim F_{p,n-p}$.
Therefore, it holds that
\begin{align*}
\bP_{\left(\mathcal{U}^L, \bm{U}\right)}\left(\bm{\beta}_0\notin \Gamma_{\alpha}^{\mathrm{coef}}(\bm{Y}^{\obs})\right)
&\le 
\bP_{\bm U}\left(T(\bm{Y}^{\obs},(\Pi_0,\bm{\beta}_0))> F_{p,n-p}^{-1}(\alpha)\right)
+
\bP_{\left(\mathcal{U}^L, \bm{U}\right)}\left(\Pi_0\notin \mC^{(L)}\right)\\
&\le (1-\alpha)+\delta_L,
\end{align*}
which is equivalent to the claimed coverage lower bound
$\bP_{\left(\mathcal{U}^L, \bm{U}\right)}\left(\bm{\beta}_0\in \Gamma_{\alpha}^{\mathrm{coef}}(\bm{Y}^{\obs})\right)\ge \alpha-\delta_L$.
\end{proof}


\subsection{Proof of Theorem \ref{thm:high-probability equivalence}}\label{subsec: proof of hungarian} Throughout this subsection, let $\bm{Y}$ be the abbreviation of $\bm{Y}^{\obs}$. Also, set
\begin{align*}
B_Y(\xi)&:= \|\Pi_0\bm X\bm\beta_0\|_2^2 + 2\sigma_0\|\Pi_0\bm X\bm\beta_0\|_2\sqrt{2\log\frac{2}{\xi}}
 + \sigma_0^2 \left(n+2\sqrt{n\log\frac{2}{\xi}}+2\log\frac{2}{\xi}\right),
\end{align*}
and
\begin{align}
    B_{\rm diag}(\xi)
&:= \|(I_n-M_{\bm X})\Pi_0\bm X\bm\beta_0\|_2^2+\sigma_0^2 
+2\sigma_0\|(I_n-M_{\bm X})\Pi_0\bm X\bm\beta_0\|_2\sqrt{2\log\frac{16n}{\xi}} \label{eq:bdiag}
\\ &+2\sigma_0^2 \left(\sqrt{\log\frac{16n}{\xi}}+\log\frac{16n}{\xi}\right). \nonumber
\end{align}

Further, define 
\begin{equation}\label{eq:eta-op-def}
\eta_{\rm op}(\xi) :=   
\frac{8 B_Y(\xi)\cdot\sqrt{k+2\sqrt{k\log \frac{4k+2}{\xi}}+2\log \frac{4k+2}{\xi}}}
{\sqrt{(n-p)-2\sqrt{(n-p)\log \frac{2}{\xi}}}-2\sqrt{k+2\sqrt{k\log \frac{4k+2}{\xi}}+2\log \frac{4k+2}{\xi}}}.
\end{equation}

Recall that the objective functions of the main optimization problem \eqref{eq:repro_Pi} and the score-weighted LAP \eqref{def:tilde_Pi} are given by
\[
F(\Pi)=\left\|(I_n-M_{(\Pi\bm X,\bm U^*)}) \bm Y\right\|_2^2,
\qquad
G(\Pi)=G_0(\Pi)+R(\Pi),
\]
where $G_0(\Pi):=\sum_i(Y_i-m_{\pi(i)})^2$ with $\bm m:=M_{(\bm X,\bm U^*)}\bm Y$ and
\[
R(\Pi)=\sum_{i=1}^n  \left\{\lambda_1   \1\{\pi(i)\neq i\}-\lambda_2   \1\{\pi(i)=i\} (Y_i-m_i)^2\right\}.
\]
The following lemma gives the condition that the solution of the score-weighted LAP is localized within $\mP_{n,k}$.

\begin{lemma}\label{lem:lap-minimizer-in-Pnk}
Assume $\lambda_1>0$ and define
\[
\tilde{\Pi}\in\argmin_{\Pi\in\mP_n} G(\Pi),
\qquad
G(\Pi)=\sum_{i=1}^n (Y_i-m_{\pi(i)})^2
+\sum_{i=1}^n\left\{\lambda_1\1\{\pi(i)\neq i\}-\lambda_2\1\{\pi(i)=i\}(Y_i-m_i)^2\right\},
\]
where $\bm m:=M_{(\bm X,\bm U^*)}\bm Y$.
Then, it holds that
\begin{equation}\label{eq:sparsity-bound-by-residual}
d(\tilde{\Pi},I_n)\ \le\ \frac{\|\bm Y-\bm m\|_2^2}{\lambda_1}.
\end{equation}
In particular, for any integer $k\ge 1$, if $\lambda_1\ge \|\bm Y-\bm m\|_2^2/k$, then it holds that $\tilde{\Pi}\in\mP_{n,k}$.
Moreover, since $\|\bm Y-\bm m\|_2^2\le \|\bm Y\|_2^2$, on the event $\{\|\bm Y\|_2^2\le B_Y(\xi)\}$ it holds that
$\tilde{\Pi}\in\mP_{n,k}$ whenever $\lambda_1\ge B_Y(\xi)/k$.
\end{lemma}

\begin{proof}
Write $r:=\|\bm Y-\bm m\|_2^2=\sum_{i=1}^n (Y_i-m_i)^2$.
For any $\Pi\in\mP_n$, we have $G_0(\Pi):=\sum_{i=1}^n (Y_i-m_{\pi(i)})^2\ge 0$ and
\[
R(\Pi)
=\lambda_1 d(\Pi,I_n)-\lambda_2\sum_{i:\pi(i)=i}(Y_i-m_i)^2
\ge \lambda_1 d(\Pi,I_n)-\lambda_2 r,
\]
since $\sum_{i:\pi(i)=i}(Y_i-m_i)^2\le r$. Therefore, we have
\[
G(\Pi)=G_0(\Pi)+R(\Pi)\ \ge\ \lambda_1 d(\Pi,I_n)-\lambda_2 r.
\]
For the identity permutation $I_n$, we have
\[
G(I_n)=\sum_{i=1}^n\{(Y_i-m_i)^2-\lambda_2(Y_i-m_i)^2\}=(1-\lambda_2)r.
\]
By its optimality, we have $G(\tilde{\Pi})\le G(I_n)$, which yields that
\[
\lambda_1 d(\tilde{\Pi},I_n)-\lambda_2 r
\le G(\tilde{\Pi})
\le (1-\lambda_2)r,
\]
This implies that $\lambda_1 d(\tilde{\Pi},I_n)\le r$ in \eqref{eq:sparsity-bound-by-residual}.
The remaining claims follow immediately, and $\|\bm Y-\bm m\|_2^2\le \|\bm Y\|_2^2$ holds because $\bm m$ is an orthogonal projection of $\bm Y$.
\end{proof}



For any edge $(\Pi,\tau)$ in the Cayley graph generated by transpositions, write the one–step increments
\[
\Delta_H^{\mathrm{edge}}(\Pi,\tau):=H(\Pi\circ\tau)-H(\Pi), \qquad H\in\{F,G_0,G\}.
\]

For $\Pi_{(0)},\cdots, \Pi_{(r)}\in\mP_{n,k}$, we work along a fixed transposition path:
\[
\Pi^{(0)}\xrightarrow{\tau_0}\Pi_{(1)}\xrightarrow{\tau_1}\cdots
\xrightarrow{\tau_{r-1}}\Pi_{(r)},
\qquad \Pi_{(\ell+1)}=\Pi_{(\ell)} \circ \tau_\ell,
\]
and denote the set of visited permutations by $\mS_{\rm path}:=\{\Pi_{(0)},\cdots,\Pi_{(r)}\}$, so that we have $|\mS_{\rm path}|=m=r+1$. When we allow all transpositions as generators, one may always take $r\le 2k$.

First, we show a path-wise comparison of edge increments, where the following bounds control the discrepancy between two objective functions $F$ and $G_0$ only along the given path.

\begin{proposition}\label{prop:edge-compare-path}
Let $\Pi_{(0)}\to\Pi_{(1)}\to\cdots\to\Pi_{(r)}$ be a shortest path in $\mathcal{G}_k$.
Let $\mS_{\rm path} = \{\Pi_{(0)},\ldots,\Pi_{(r)}\}$ and $m=r+1$.
From Proposition~\ref{prop:finite-path-op}, on the event
\begin{align*}
\mathcal{E}_{\rm path}(\xi)
&:= \Bigg\{
\max_{\Pi'\in\mS_{\rm path}}
\|\Pi' M_{(\bm X,\bm U^*)}-M_{(\Pi'\bm X,\bm U^*)}\|_{\rm op}\\
&\le 
\frac{4\sqrt{k+2\sqrt{k\log\frac{4k+2}{\xi}}+2\log\frac{4k+2}{\xi}}}
{\sqrt{(n-p)-2\sqrt{(n-p)\log \frac{2}{\xi}}}-2\sqrt{k+2\sqrt{k\log \frac{4k+2}{\xi}}+2\log \frac{4k+2}{\xi}}}
\Bigg\},
\end{align*}
for every edge $(\Pi_{(\ell)},\tau_\ell)$ along the path, it holds that
\begin{align*}
&\left|\frac{1}{2} \Delta_{G_0}^{\mathrm{edge}}(\Pi_{(\ell)},\tau_\ell)-\Delta_F^{\mathrm{edge}}(\Pi_{(\ell)},\tau_\ell)\right|\\
&\le  
\frac{8 \|\bm Y\|_2^2 \sqrt{k+2\sqrt{k\log\frac{4k+2}{\xi}}+2\log\frac{4k+2}{\xi}}}
{\sqrt{ (n-p)-2\sqrt{(n-p)\log(2/\xi)}}-2\sqrt{k+2\sqrt{k\log\frac{4k+2}{\xi}}+2\log\frac{4k+2}{\xi}}}.
\end{align*}

Consequently, for the same edge, we have
\begin{align*}
    \left| \Delta_G^{\mathrm{edge}}(\Pi_{(\ell)},\tau_\ell)-\Delta_{G_0}^{\mathrm{edge}}(\Pi_{(\ell)},\tau_\ell)\right|
    =\left|\Delta_R^{\mathrm{edge}}(\Pi_{(\ell)},\tau_\ell)\right|
    \le 2\lambda_1+2\lambda_2\max_{i\in[n]} (Y_i-m_i)^2.
\end{align*}
\end{proposition}


\begin{proof}
By the identity
\begin{align}    
&\frac{1}{2}\Delta_{G_0}^{\mathrm{edge}}-\Delta_F^{\mathrm{edge}}\\
&=\bm Y^\top  \left(\Pi_{(\ell)}M_{(\bm X,\bm U^*)}-M_{(\Pi_{(\ell)}\bm X,\bm U^*)}\right)  \bm Y
-\bm Y^\top  \left( (\Pi_{(\ell)} \circ \tau_\ell) M_{(\bm X,\bm U^*)}-M_{((\Pi_{(\ell)} \circ \tau_\ell)\bm X,\bm U^*)}\right)  \bm Y, 
\end{align}
the triangle inequality and Proposition \ref{prop:finite-path-op} give the first inequality. 

Furthermore, only two rows are affected by an adjacent transposition $\tau_{(\ell)}$, hence the $\lambda_1$–term changes by at most $2\lambda_1$. The $\lambda_2$–term involves at most two diagonal positions and contributes at most $2\lambda_2\max_i (Y_i-m_i)^2$. This gives the second inequality.
\end{proof}

Next step is to derive an explicit relation between $C_{\min}(\Pi_0)$ and $\min_{\Pi\in\mP_{n,k}, \Pi\neq\hat\Pi}\left\{F(\Pi)-F(\hat\Pi)\right\}$.

We also define the corresponding (random) separation gap
\begin{equation}\label{eq:delta-F}
\Delta_F
:= \min_{\Pi\in\mP_{n,k}\setminus\{\hat\Pi\}}
\left\{F(\Pi)-F(\hat\Pi)\right\},
\end{equation}
where $\hat\Pi=\hat\Pi(\bm U,\bm U^*)$ is the minimizer of $F(\cdot)$ over $\mP_{n,k}$.
We adopt the convention that $\Delta_F=0$ if the minimizer of $F$ over $\mP_{n,k}$ is not unique.


\begin{proposition}\label{prop:F separation gap}
Assume $n-p \ge 2$ and recall $m_k := |\mP_{n,k}|$.
Define the functions
\begin{equation}\label{eq:theta1}
\theta_1(\xi) := \left(\frac{\xi}{6m_k}\right)^{\frac{1}{n-p-1}},
\qquad
\gamma_1(\xi) := \cos\left(\theta_1(\xi)\right),
\end{equation}
and
\begin{equation}\label{eq:bar-BU}
\bar B_U(\xi):=1+2\sqrt{\frac{\log(6/\xi)}{n}}+2\frac{\log(6/\xi)}{n}.
\end{equation}
Furthermore, define
\begin{equation}\label{eq:Delta-under-xi}
\underline\Delta_\xi :=
(1-\gamma_1(\xi)^2)C_{\min}
-2\sigma_0\sqrt{C_{\min}\bar B_U(\xi)}
-\sigma_0^2 \bar B_U(\xi).
\end{equation}
Then, with probability at least $1-\xi$, the separation gap $\Delta_F$ satisfies
\[
\Delta_F \ge n\underline\Delta_\xi.
\]
In particular, if $\underline\Delta_\xi>0$, then we have $\hat\Pi=\Pi_0$ and $\Delta_F>0$.
\end{proposition}

\begin{proof}
Define $A_\Pi := M_{(\Pi\bm X,\bm U^*)}$ and $A_0:=M_{(\Pi_0\bm X,\bm U^*)}$.
Fix $\Pi\in\mP_{n,k}\setminus\{\Pi_0\}$.
Since it holds that $\Pi_0\bm X\bm\beta_0\in \mathrm{range}(A_0)$, we have
\begin{align*}
F(\Pi)-F(\Pi_0)
&=\|(I_n-A_\Pi)\Pi_0\bm X\bm\beta_0\|_2^2
+2\sigma_0\bm U^\top (I_n-A_\Pi)\Pi_0\bm X\bm\beta_0
-\sigma_0^2 \bm U^\top (A_\Pi-A_0)\bm U.
\end{align*}
Define $\bm w_\Pi:=(I_n-M_{\Pi\bm X})\Pi_0\bm X\bm\beta_0$.
A direct calculation gives
\[
\|(I_n-A_\Pi)\Pi_0\bm X\bm\beta_0\|_2^2
=\|\bm w_\Pi\|_2^2\left(1-\phi\left((I_n-M_{\Pi\bm X})\bm U^*,\bm w_\Pi\right)\right).
\]
Define the events
\[
\mathcal{E}_1(\xi):=\left\{\max_{\Pi\in\mP_{n,k}}
\phi\left((I_n-M_{\Pi\bm X})\bm U^*,\bm w_\Pi\right)\le \gamma_1(\xi)^2\right\},
\qquad
\mathcal{E}_3(\xi):=\left\{\|\bm U\|_2^2\le n\bar B_U(\xi)\right\}.
\]
By Lemma~\ref{lem:max-sim} and the choice of $\theta_1(\xi)$, we have
$\bP(\mathcal{E}_1(\xi)^c)\le \xi/3$.
By a standard $\chi^2$ upper-tail bound, we have
$\bP(\mathcal{E}_3(\xi)^c)\le \xi/3$.
Hence $\bP(\mathcal{E}_1(\xi)\cap \mathcal{E}_3(\xi))\ge 1-2\xi/3\ge 1-\xi$.

Furthermore, on $\mathcal{E}_1(\xi)$, we have
\[
\|(I_n-A_\Pi)\Pi_0\bm X\bm\beta_0\|_2^2\ge (1-\gamma_1(\xi)^2)\|\bm w_\Pi\|_2^2.
\]
For the cross term, using $\|(I_n-A_\Pi)\bm U\|_2\le \|\bm U\|_2$, it yields that
\[
\left|2\sigma_0\bm U^\top (I_n-A_\Pi)\Pi_0\bm X\bm\beta_0\right|
\le 2\sigma_0\|\bm U\|_2\|\bm w_\Pi\|_2.
\]
For the quadratic term, since $A_\Pi-A_0$ is a difference of orthogonal projections, we have
\[
\sigma_0^2 \bm U^\top (A_\Pi-A_0)\bm U
\le \sigma_0^2 \|\bm U\|_2^2.
\]
Therefore, on $\mathcal{E}_1(\xi)\cap \mathcal{E}_3(\xi)$, we have
\[
F(\Pi)-F(\Pi_0)
\ge (1-\gamma_1(\xi)^2)\|\bm w_\Pi\|_2^2
-2\sigma_0\sqrt{n\bar B_U(\xi)}\|\bm w_\Pi\|_2
-\sigma_0^2 n\bar B_U(\xi).
\]
Using the definition of $C_{\min}$, we obtain
$\|\bm w_\Pi\|_2^2\ge n C_{\min}\cdot \max\{d(\Pi_0,I_n)-d(\Pi,I_n),1\}\ge nC_{\min}$,
and a direct calculation yields
$F(\Pi)-F(\Pi_0)\ge n\underline\Delta_\xi$.
Taking the minimum over $\Pi\neq \Pi_0$ gives $\Delta_F\ge n\underline\Delta_\xi$.
If $\underline\Delta_\xi>0$, then it holds that $F(\Pi)>F(\Pi_0)$ for all $\Pi\neq \Pi_0$, hence the claim that $\hat\Pi=\Pi_0$ holds.
\end{proof}




We are now in position to prove the main result (Theorem \ref{thm:high-probability equivalence}). 

\begin{proof}
Remind that 
\begin{align*}
B_Y(\xi) & = \|\Pi_0\bm X\bm\beta_0\|_2^2 + 2\sigma_0\|\Pi_0\bm X\bm\beta_0\|_2\sqrt{2\log\frac{2}{\xi}}
 + \sigma_0^2 \left(n+2\sqrt{n\log\frac{2}{\xi}}+2\log\frac{2}{\xi}\right),
\end{align*}
and
\begin{align*}
    B_{\rm diag}(\xi)
 = &\|(I_n-M_{\bm X})\Pi_0\bm X\bm\beta_0\|_2^2+\sigma_0^2
+2\sigma_0\|(I_n-M_{\bm X})\Pi_0\bm X\bm\beta_0\|_2\sqrt{2\log\frac{16n}{\xi}}
\\ 
&+2\sigma_0^2 \left(\sqrt{\log\frac{16n}{\xi}}+\log\frac{16n}{\xi}\right).
\end{align*}
By Lemma \ref{lem:Y-and-diag}, with probability at least $1-\xi$, we also have
$\|\bm Y\|_2^2\le B_Y(\xi)$.
Define the event
\[
\mathcal{E}_{\rm op}(\xi):=\mathcal E_{\rm path}(\xi)\cap\{\|\bm Y\|_2^2\le B_Y(\xi)\}.
\]
On $\mathcal{E}_{\rm op}(\xi)$, Proposition \ref{prop:edge-compare-path} yields the uniform edge-wise control:
\begin{equation}\label{eq:edge-eta}
\left|
\frac{1}{2}\Delta_{G_0}^{\rm edge}(\Pi_{(\ell)},\tau_\ell)-\Delta_F^{\rm edge}(\Pi_{(\ell)},\tau_\ell)
\right|
\le\eta_{\rm op}(\xi),
\qquad \ell=0,\ldots,r-1,
\end{equation}
with $\eta_{\rm op}(\xi)$ as in \eqref{eq:eta-op-def}.

For the regularizer
$
R(\Pi)=\lambda_1\sum_{i}\1\{\pi(i)\neq i\}-\lambda_2\sum_{i}\1\{\pi(i)=i\}(Y_i-m_i)^2,
$
a single transposition affects at most two rows, hence
\begin{equation}\label{eq:edge-R}
\left|\Delta_R^{\rm edge}(\Pi_{(\ell)},\tau_\ell)\right| \le 2\lambda_1 + 2\lambda_2 \max_{1\le i\le n}(Y_i-m_i)^2.
\end{equation}
By Lemma~\ref{lem:Y-and-diag} (in particular \eqref{eq:diag-sq-finite}), with probability at least $1-\xi$, we have
\[
\max_{1\le i\le n} (Y_i-m_i)^2
\le \max_{1\le i\le n}\mathbb{E}\big[(Y_i-m_i)^2\mid \bm U^*\big]
     + \max_{1\le i\le n}\Big|(Y_i-m_i)^2-\mathbb{E}\big[(Y_i-m_i)^2\mid \bm U^*\big]\Big|.
\]
Moreover, writing $\bm\mu:=(I_n-M_{(X,\bm U^*)})\Pi_0X\beta_0$, we have
$\mathbb{E}\big[(Y_i-m_i)^2\mid \bm U^*\big]=\mu_i^2+\sigma_0^2\|(I_n-M_{(X,\bm U^*)})e_i\|_2^2\le \|\bm\mu\|_\infty^2+\sigma_0^2$,
and $\|\bm\mu\|_\infty\le \|\bm\mu\|_2\le \|(I_n-M_X)\Pi_0X\beta_0\|_2$.
Set
\[
\mathcal{E}_{\rm diag}(\xi):=\left\{\max_{1\le i\le n}(Y_i-m_i)^2\le B_{\rm diag}(\xi)\right\}.
\]
On $\mathcal{E}_{\rm diag}(\xi)$, \eqref{eq:edge-R} gives that, for every edge along the path,
\begin{equation}\label{eq:edge-R-ub}
\left|\Delta_R^{\rm edge}(\Pi_{(\ell)},\tau_\ell)\right| \le 2\lambda_1+2\lambda_2 B_{\rm diag}(\xi),
\qquad \ell=0,\ldots,r-1.
\end{equation}

Summing \eqref{eq:edge-eta} over the $r\le 2k$ edges and using
$
H(\Pi)-H(\hat \Pi)=\sum_{\ell=0}^{r-1}\Delta_H^{\rm edge}(\Pi_{(\ell)},\tau_\ell)
$
for $H\in\{F,G_0\}$ yields
\begin{equation}\label{eq:G0-vs-F}
\left|
\frac{1}{2}\left(G_0(\Pi)-G_0(\hat\Pi)\right)-\left(F(\Pi)-F(\hat\Pi)\right)
\right|
 \le r \eta_{\rm op}(\xi) \le 2k \eta_{\rm op}(\xi).
\end{equation}
Likewise, summing \eqref{eq:edge-R-ub} along the path, we have
\begin{equation}\label{eq:G-vs-G0}
\left|\left(G(\Pi)-G(\hat\Pi)\right)-\left(G_0(\Pi)-G_0(\hat\Pi)\right)\right|
 \le r\left(2\lambda_1+2\lambda_2 B_{\rm diag}(\xi)\right) \le 2k\left(2\lambda_1+2\lambda_2 B_{\rm diag}(\xi)\right).
\end{equation}
Combining them yields the following inequalities:
\begin{align}
G(\Pi)-G(\hat\Pi)
& \ge 2\left(F(\Pi)-F(\hat\Pi)\right) - 2r \eta_{\rm op}(\xi) - r \left(2\lambda_1+2\lambda_2 B_{\rm diag}(\xi)\right)\nonumber\\
& \ge 2\Delta_F - 4k \eta_{\rm op}(\xi) - 4k\lambda_1 - 4k\lambda_2 B_{\rm diag}(\xi).
\label{eq:G-lower-final}
\end{align}

By Proposition \ref{prop:F separation gap}, the following event
\[
\mathcal{E}_{\rm gap}(\xi) := \left\{\min_{\Pi\in\mP_{n,k}, \Pi \neq \hat\Pi}\left\{F(\Pi)-F(\hat {\Pi}) \right\} \ge n\underline\Delta_\xi \right\},
\] 
has a probability at least $1 - \xi$.

On the intersection event
\[
\mathcal{E}_{\rm op}(\xi)\cap\mathcal{E}_{\rm diag}(\xi)\cap\mathcal{E}_{\rm gap}(\xi),
\]
the inequality \eqref{eq:G-lower-final} and the condition \eqref{eq:lambda-window} imply that, for every $\Pi\in\mP_{n,k}\setminus\{\hat\Pi\}$, we have
\[
G(\Pi)-G(\hat\Pi) \ge 2\Delta_F - 4k \eta_{\rm op}(\xi) - 4k\lambda_1 - 4k\lambda_2 B_{\rm diag}(\xi)
 > 0.
\]
Therefore, $G$ has a unique minimizer over $\mP_{n,k}$ and it must equal $\hat\Pi$.
Finally, by Lemmas~\ref{lem:perm-diff-finite} and~\ref{lem:denom-lb} (applied with $\xi/2$) and Lemma~\ref{lem:Y-and-diag}, we have
\[
\bP(\mathcal{E}_{\rm op}(\xi))\ge 1-2\xi,\qquad
\bP(\mathcal{E}_{\rm diag}(\xi))\ge 1-\xi,\qquad
\bP(\mathcal{E}_{\rm gap}(\xi))\ge 1-\xi,
\]
so a union bound yields the desired probability lower bound. By Lemma~\ref{lem:lap-minimizer-in-Pnk}, on the event $\{\|\bm Y\|_2^2\le B_Y(\xi)\}$ and if $\lambda_1\ge B_Y(\xi)/k$,
every minimizer of $G$ over $\mP_n$ belongs to $\mP_{n,k}$.
Hence, we obtain $\tilde\Pi\in\mP_{n,k}$ and therefore $\tilde\Pi=\hat\Pi$ by the uniqueness of the minimizer over $\mP_{n,k}$.
\end{proof}

\section{Additional theoretical results}\label{section:appendix theory}

\subsection{A counterexample when $n-2k < p$}
\begin{theorem}\label{thm:counterexample}
    Let $p<n$ and $n-2k \le p - 1$. Further, assume $k \ge 2$ so that a nontrivial transposition is permitted.
Then, there exist a full-column-rank matrix $\bm{X} \in \R^{n\times p}$, a pair of distinct permutations $\Pi_0,\Pi_1 \in \mP_{n,k}$,
and coefficient vectors $\bm{\beta}_0 \ne \bm{\beta}_1 \in \R^p$ such that
\[
\Pi_0 \bm{X} \bm{\beta}_0 = \Pi_1 \bm{X} \bm{\beta}_1
\]
holds. Consequently, the minimizers of \eqref{eq:repro_ideal} are not unique in general, even with $\sigma_0 = 0$.
\end{theorem}

\begin{proof}
Fix two distinct indices $a\neq b$ in $[n]$. Set $\Pi_0:=I_n$ and let $\Pi_1$ be the transposition that swaps $a$ and $b$.
Then $d(\Pi_1,I_n)=2\le k$ because $k\ge2$, hence we have $\Pi_0,\Pi_1\in\mP_{n,k}$.

Set
\[
t:=\max\{n-2k,0\}.
\]
By the assumption $n-2k\le p-1$ we have $t\le p-1$. Choose any subset $T\subset[n]$ with $|T|=t$ and set $S:=[n]\setminus T$.

We construct a full-column-rank design $\bm X\in\R^{n\times p}$ together with $\bm\beta_0\ne\bm\beta_1$
and $\Pi_0\bm X\bm\beta_0=\Pi_1\bm X\bm\beta_1$.

First, consider the case $p=1$. Put $\bm X\in\R^{n\times1}$ with $X_a=1$, $X_b=-1$, and $X_i=0$ for $i\notin\{a,b\}$.
Then, we have $\rank(\bm X)=1$ (full column rank). Take $\beta_0=1$ and $\beta_1=-1$.
We have $(\bm X\beta_0)_a=1$ and $(\bm X\beta_0)_b=-1$, while $\bm X\beta_1$ equals $-1$ at $a$ and $1$ at $b$,
so applying the transposition $\Pi_1$ swaps these two values and yields
\[
\Pi_1\bm X\beta_1=\bm X\beta_0=\Pi_0\bm X\beta_0,
\]
with $\beta_0\ne\beta_1$.

Next, consider the case $p\ge2$.
Let $v:=e_1\in\R^p$ and $w:=e_2\in\R^p$, so $w\in v^\perp$ and $\|v\|_2^2=1$. Define two row vectors
\[
r_a:=w+v,\qquad r_b:=w-v,
\]
so that we have $r_a v=1$ and $r_b v=-1$.

We define the rows of $\bm X$ as follows.
\begin{itemize}
\item Put $X_a:=r_a$ and $X_b:=r_b$.
\item For the $t$ indices in $T$, assign distinct standard basis vectors among $\{e_2,e_3,\dots,e_{t+1}\}$:
for example, after ordering $T=\{i_1,\dots,i_t\}$, set $X_{i_j}:=e_{j+1}$.
This gives $\rank(\bm X_T)=t$ and $\bm X_T v=0$.
\item For the remaining indices in $S\setminus\{a,b\}$, assign the missing basis vectors among $\{e_2,\dots,e_p\}$
so that the set of rows contains $\{e_2,\dots,e_p\}$ at least once; if there are extra rows, assign any vectors in $v^\perp$. This is feasible because $n\ge p+1$ (since $p<n$), hence
$|S\setminus\{a,b\}| = n-t-2 \ge p-t-1$.
\end{itemize}
By construction, all rows except $a,b$ lie in $v^\perp$, hence we have $X_i v=0$ for all $i\notin\{a,b\}$.
Moreover, the row set contains $\{r_a,e_2,\dots,e_p\}$, which is linearly independent, so $\rank(\bm X)=p$
(full column rank) holds.

Now choose $\bm\beta_0:=-\tfrac12 v+\tfrac12 w$ and $\bm\beta_1:=\bm\beta_0+v=\tfrac12 v+\tfrac12 w$.
Then, we have
\[
r_a\bm\beta_0=(v+w)\cdot\left(-\tfrac12 v+\tfrac12 w\right)=0,\qquad
r_b\bm\beta_0=(w-v)\cdot\left(-\tfrac12 v+\tfrac12 w\right)=1,
\]
and
\[
r_a\bm\beta_1=(v+w)\cdot\left(\tfrac12 v+\tfrac12 w\right)=1,\qquad
r_b\bm\beta_1=(w-v)\cdot\left(\tfrac12 v+\tfrac12 w\right)=0.
\]
For every $i\notin\{a,b\}$, because of $X_i\in v^\perp$ we have $X_i(\bm\beta_1-\bm\beta_0)=X_i v=0$, hence it holds that $X_i\bm\beta_1=X_i\bm\beta_0$.
Therefore, $\bm X\bm\beta_1$ equals $\bm X\bm\beta_0$ at all rows except $a,b$, where the two values are swapped.
Applying the transposition $\Pi_1$ thus yields
\[
\Pi_1\bm X\bm\beta_1=\bm X\bm\beta_0=\Pi_0\bm X\bm\beta_0,
\]
with $\bm\beta_0\ne\bm\beta_1$.

In either case, we construct the design matrix $\bm X$ with full column-rank, the permutations $\Pi_0\ne\Pi_1$ in $\mP_{n,k}$, and the coefficients $\bm\beta_0\ne\bm\beta_1$
satisfying $\Pi_0\bm X\bm\beta_0=\Pi_1\bm X\bm\beta_1$. This yilelds the non-uniqueness of \eqref{eq:repro_ideal} even when $\sigma_0=0$.
\end{proof}


\subsection{Theory on the partially permuted designs}

\begin{proposition}\label{prop:identification partial}
    Assume that the rows of design matrix $\bm{X}$ and $\bm{Z}$ are i.i.d. generated by a continuous distribution. Also, suppose $n - 2k \ge p_1 + p_2$. Then, true parameters
$(\Pi_0,\bm\beta_{0,1},\bm\beta_{0,2},\sigma_0)$ can be uniquely recovered by the following optimization problem:
\begin{equation}\label{eq:repro-ideal-partial}
(\Pi_0,\bm\beta_{0,1},\bm\beta_{0,2},\sigma_0)
=
\argmin_{\Pi\in\mP_{n,k},\ \bm\beta_1\in\R^{p_1},\ \bm\beta_2\in\R^{p_2},\ \sigma\ge 0}
\left\|
\bm Y^{\obs}-\Pi \bm X\bm\beta_1-\bm Z\bm\beta_2-\sigma \bm u^{\rel}
\right\|^2 .
\end{equation}
For a fixed $\Pi \in \mP_n$, define an expanded design matrix $W(\Pi) = [\Pi \bm{X}, \bm{Z}]$. Furthermore, it holds that
\begin{equation}
   \Pi_0 = \argmin_{\Pi \in \mP_{n,k}}\| (I_n - M_{(W(\Pi), \bm{u}^{\rel})}) \bm{Y}^{\obs}\|_2^2,
\end{equation}
where $M_{(W(\Pi), \bm{u}^{\rel})}$ is the projection matrix to the range of $(W(\Pi), \bm{u}^{\rel})$. 
\end{proposition}

\begin{proof}
Applying same strategy in the proof of Proposition \ref{prop:identification} and treating $M_{(\Pi \bm{X}, \bm{u}^{\rel})}$ as $M_{(W(\Pi), \bm{u}^{\rel})}$ yield the desired result.
\end{proof}

Recall the partially permuted model \eqref{eq:model_partial} and write $p:=p_1+p_2$.
For any $\Pi\in\mP_{n,k}$ define $W(\Pi):=[\Pi\bm X,\bm Z]$ and $M_{W(\Pi)}$ as in
\eqref{eq:Wpartial-def}.  The corresponding separation constant is given by
\begin{equation}\label{eq:c-min-partial}
C_{\min}^{\mathrm{partial}}(\Pi_0)
:=\min_{\substack{\Pi\in\mP_{n,k}\\ \Pi\neq \Pi_0}}
\frac{\left\| (I_n-M_{W(\Pi)})W(\Pi_0)\binom{\bm\beta_{1,0}}{\bm\beta_{2,0}}\right\|_2^2}
     {n\max\{d(\Pi_0,I_n)-d(\Pi,I_n),1\}}
=\min_{\substack{\Pi\in\mP_{n,k}\\ \Pi\neq \Pi_0}}
\frac{\left\| (I_n-M_{(\Pi\bm X,\bm Z)})\Pi_0\bm X\bm\beta_{1,0}\right\|_2^2}
     {n\max\{d(\Pi_0,I_n)-d(\Pi,I_n),1\}}.
\end{equation}


\begin{theorem}\label{thm:recovery partial}
Suppose $C_{\min}^{\mathrm{partial}}(\Pi_0) > 0$ and $n > p_1 + p_2 + 1$ hold.
Let $\Delta_{\mathrm{partial}}(\gamma)$ be defined as in~\eqref{eq:Delta-gamma} but with
$(C_{\min}(\Pi_0),p)$ replaced by $(C_{\min}^{\mathrm{partial}}(\Pi_0),p_1+p_2)$.

Then, for any $\gamma \in (0, 1/4]$ and any $L \ge 2$, we have
\begin{equation}\label{eq:recovery-finite-L-partial}
    \bP_{(\mathcal{U}^L, \bm{U})}\left( \Pi_0 \notin \mathcal{C}_{\mathrm{partial}}^{(L)} \right)
    \le
    \Delta_{\mathrm{partial}}(\gamma)
    + \left(1 - \frac{\gamma^{n-1}}{n-1}\right)^L.
\end{equation}
In particular, taking $\gamma = \gamma_L$ as in~\eqref{eq:gamma-L} gives
\begin{equation}\label{eq:recovery-finite-L-partial-gammaL}
    \bP_{(\mathcal{U}^L, \bm{U})}\left( \Pi_0 \notin \mathcal{C}_{\mathrm{partial}}^{(L)} \right)
    \le
    \delta_L^{\mathrm{partial}}
    :=
    \Delta_{\mathrm{partial}}(\gamma_L)
    + \left(1 - \frac{\gamma_L^{n-1}}{n-1}\right)^L
    = O\left( \frac{(\log L)^{(n-(p_1+p_2))/2}}{L^{(n-(p_1+p_2)-1)/(2(n-1))}} \right).
\end{equation}
\end{theorem}

\begin{proof}
This is an immediate application of Theorem~\ref{thm:recovery} to the design matrice
$\bm{W}(\Pi) = (\Pi \bm{X}, \bm{Z})$ with its total dimension $p=p_1+p_2$, for which the separation constant is
$C_{\min}^{\mathrm{partial}}(\Pi_0)$ as defined above.
\end{proof}


\begin{theorem}\label{thm:p-value partial}
    Let $\bm{U}^*$ be an independent copy of $\bm{U} \sim \Normal(0,I_n)$. For any $\alpha \in (0,1)$, it follows that
\begin{equation}
    \bP \left(\tilde{T}^{\mathrm{partial}}(\bm{U}, \bm{\theta}_{\mathrm{null}}) \in  \left( B_{1-\alpha}^{\mathrm{partial}}(\bm{\theta}_{\mathrm{null}}, S_{\bm{\theta}_{\mathrm{null}}}(\bm{U})) \right)^c \mid S_{\bm{\theta}_{\mathrm{null}}}^{\mathrm{partial}}(\bm{U}) = \bm{s} \right)  \le \alpha,
\end{equation}
for $\bm{\theta}_{\mathrm{null}} := (I_n, \bm{\beta}_1, \bm{\beta}_2, \sigma)$. Moreover, marginally, we have
\begin{equation}\label{eq:marginal validity}
\bP \left(\tilde{T}^{\mathrm{partial}}(\bm{U}, \bm{\theta}_{\mathrm{null}}) \in  \left( B_{1-\alpha}^{\mathrm{partial}}(\bm{\theta}_{\mathrm{null}}, S_{\bm{\theta}_{\mathrm{null}}}(\bm{U})) \right)^c\right)  \le \alpha.
\end{equation}
\end{theorem}

\begin{proof}
    Define 
    \begin{equation}\label{eq:block-proj}
            M_{W(\Pi)}-M_{\bm Z}=M_{R_{\bm Z}\Pi\bm X}=:M_{\bm X_{\perp}(\Pi)},
\qquad
\text{and}\qquad
R_{\bm Z}=M_{\bm X_{\perp}(I_n)}+R_{W(I_n)},
    \end{equation}
    where the two projectors on the right-hand side of the second equality are orthogonal. Under the null parameter $\bm\theta_{\mathrm{null}}=(I_n,\bm\beta_1,\bm\beta_2,\sigma)$, we have
    \begin{equation}\label{eq:S-partial-correct}
S^{\mathrm{partial}}_{\bm\theta_{\mathrm{null}}}(\bm U)
:=\left(M_{W_0}\bm Y_{\bm\theta_{\mathrm{null}}}, \|R_{W_0}\bm Y_{\bm\theta_{\mathrm{null}}}\|_2\right)
=\left(M_{W_0}\bm Y_{\bm\theta_{\mathrm{null}}}, \sigma \|R_{W_0}\bm U\|_2\right).
\end{equation} is (minimal)
sufficient for $(\bm\beta_1,\bm\beta_2,\sigma)$, so that the direction
\[
\bm A(\bm U):=\frac{R_{W_0}\bm U}{\|R_{W_0}\bm U\|}
\]
is ancillary. By Basu's theorem, conditionally on $S^{\mathrm{partial}}_{\bm\theta_{\mathrm{null}}}$, $\bm A(\bm U)$ is uniformly
distributed on the unit sphere of the subspace $\range(R_{W_0})$ and is independent of
$(\bm\beta_1,\bm\beta_2,\sigma)$.

Consequently, for any realization $S=(s_1,s_2)$ of $S^{\mathrm{partial}}_{\bm\theta_{\mathrm{null}}}(\bm U)$, we have the conditional representation 
\begin{equation}\label{eq:cond-rep}
\bm Y_{\bm\theta_{\mathrm{null}}}=s_1+s_2\bm A(\bm U),\qquad
R_{\bm Z}\bm Y_{\bm\theta_{\mathrm{null}}}
=M_{\bm X_{\perp}(I_n)}s_1+s_2\bm A(\bm U),
\end{equation}
where $R_{\bm Z}=M_{\bm X_{\perp}(I_n)}+R_{W_0}$ from \eqref{eq:block-proj} and
$R_{W_0}s_1=0$.

By definition, we have
\[
\tilde{T}^{\mathrm{partial}}(\bm U,\bm\theta)
:=d\left(\argmin_{\Pi\in\mP_{n,k}}\|(I_n-M_{W(\Pi)}) \bm Y_{\bm\theta}\|_2^2,\ I_n\right)
=d \left(\argmax_{\Pi\in\mP_{n,k}}\|(M_{W(\Pi)}-M_{\bm Z}) \bm Y_{\bm\theta}\|_2^2, I_n\right),
\]
where the equivalence follows because $\|M_{W(\Pi)}\bm Y\|_2^2=\|M_{\bm Z}\bm Y\|_2^2+\|(M_{W(\Pi)}-M_{\bm Z})\bm Y\|_2^2$ holds
and the first term is $\Pi$-free. Using the identity \eqref{eq:block-proj}, we have
\[
(M_{W(\Pi)}-M_{\bm Z})\bm Y_{\bm\theta_{\mathrm{null}}}
=M_{\bm X_{\perp}(\Pi)} R_{\bm Z}\bm Y_{\bm\theta_{\mathrm{null}}}
=M_{\bm X_{\perp}(\Pi)}\left(M_{\bm X_{\perp}(I_n)} s_1+s_2 \bm A(\bm U) \right),
\]
hence, conditionally on $S = (s_1, s_2)$, the only randomness in $\tilde T^{\mathrm{partial}}(\bm U,\bm\theta_{\mathrm{null}})$ stems from
$\bm A(\bm U)$, while $M_{\bm X_{\perp}(I_n)}s_1$ and $s_2$ are fixed.

Let $\bm U^*$ be an independent copy of $\bm U$ and define
$\bm A(\bm U^*):=R_{W_0}\bm U^*/\|R_{W_0}\bm U^*\|$. By the property of the ancillary statistic,
$\bm A(\bm U) \mid S  \deq \bm A(\bm U^*) \mid S $, hence it holds that
\begin{equation}\label{eq:law-equality}
\tilde T^{\mathrm{partial}}(\bm U,\bm\theta_{\mathrm{null}}) \mid  S \deq 
\tilde T^{\mathrm{partial}}(\bm U^*,\bm\theta_{\mathrm{null}}) \mid  S.
\end{equation}
Therefore, by construction of $B_{1-\alpha}^{\mathrm{partial}}(\bm\theta_{\mathrm{null}},\bm s)$, we have
\[
\bP\left(
\tilde{T}^{\mathrm{partial}}(\bm U,\bm\theta_{\mathrm{null}})
\in \left(B_{1-\alpha}^{\mathrm{partial}}(\bm\theta_{\mathrm{null}},\bm s)\right)^c
\mid S_{\bm\theta_{\mathrm{null}}}^{\mathrm{partial}}(\bm U)=\bm s
\right)\le \alpha .
\]

Taking expectation with $S$, we also have the marginal result.
\end{proof}

\begin{theorem}[Joint confidence set for $(\bm{\beta}_1,\bm{\beta}_2)$]\label{thm:confidence-set-partial-joint}
Fix any $\alpha\in(0,1)$ and consider the partially permuted model~\eqref{eq:model_partial}.
Under the conditions of Theorem~\ref{thm:recovery partial}, the confidence set
$\Gamma_{\alpha}^{\mathrm{joint}}(\bm Y^{\obs})$ defined in~\eqref{eq:confset-partial-joint}
satisfies, for any $L\ge 2$,
\[
\bP_{(\mathcal U^L,\bm U)}\left(
(\bm\beta_{1,0},\bm\beta_{2,0}) \in \Gamma_{\alpha}^{\mathrm{joint}}(\bm Y^{\obs})
\right)
\ge \alpha-\delta_L^{\mathrm{partial}},
\]
where $\delta_L^{\mathrm{partial}}$ is defined in~\eqref{eq:recovery-finite-L-partial-gammaL}.
\end{theorem}

\begin{proof}
Set $p:=p_1+p_2$.
Recall that
\begin{equation}\label{eq:confset-partial-joint}
\Gamma_\alpha^{\mathrm{joint}}(\bm Y^{\obs})
:=
\bigcup_{\Pi \in \mathcal{C}_{\mathrm{partial}}^{(L)}(\bm Y^{\obs})}
\left\{
\bm\beta \in \mathbb R^{p}:
T_{\mathrm{partial}}(\bm Y^{\obs},\Pi,\bm\beta)
\le F^{-1}_{p,n-p}(\alpha)
\right\},
\end{equation}
where the statistic $T(\bm Y^{\obs},(\Pi,\bm\beta_1,\bm\beta_2))$ is defined in~\eqref{eq:Tjointpartial-def}.

On $\{\Pi_0\in\mathcal C_{\mathrm{partial}}^{(L)}(\bm Y^{\obs})\}$, we have
\[
(\bm\beta_{1,0},\bm\beta_{2,0})\notin\Gamma_{\alpha}^{\mathrm{joint}}(\bm Y^{\obs})
\ \Longrightarrow\
T(\bm Y^{\obs},(\Pi_0,\bm\beta_{1,0},\bm\beta_{2,0}))>F^{-1}_{p,n-p}(\alpha).
\]
Hence, by Theorem~\ref{thm:recovery partial}, it holds that
\begin{align*}
\bP\left((\bm\beta_{1,0},\bm\beta_{2,0})\notin\Gamma_{\alpha}^{\mathrm{joint}}(\bm Y^{\obs})\right)
&\le \delta_L^{\mathrm{partial}}
+\bP\left(T(\bm Y^{\obs},(\Pi_0,\bm\beta_{1,0},\bm\beta_{2,0}))>F^{-1}_{p,n-p}(\alpha)\right).
\end{align*}

For the second term, under \eqref{eq:model_partial} we have
\[
\bm Y^{\obs}-\Pi_0\bm X\bm\beta_{1,0}-\bm Z\bm\beta_{2,0}
=\sigma_0\bm U,
\qquad \bm U\sim N(0,I_n).
\]
Define $M_0:=M_{(\Pi_0\bm X,\bm Z)}$ and $P_0:=I_n-M_0$. Under the rank condition
$\rank(\Pi_0\bm X,\bm Z)=p$, $M_0\bm U$ and $P_0\bm U$ are independent Gaussian vectors with
\[
\|M_0\bm U\|_2^2\sim \chi^2_p,
\qquad
\|P_0\bm U\|_2^2\sim \chi^2_{n-p},
\]
and they are independent. Therefore, we have
\[
T(\bm Y^{\obs},(\Pi_0,\bm\beta_{1,0},\bm\beta_{2,0}))
=\frac{\|M_0\bm U\|_2^2/p}{\|P_0\bm U\|_2^2/(n-p)}
\sim F_{p,n-p},
\]
which implies $\bP\left(T(\bm Y^{\obs},(\Pi_0,\bm\beta_{1,0},\bm\beta_{2,0}))\le F^{-1}_{p,n-p}(\alpha)\right)=\alpha$. Combining the bounds yields
\[
\bP\left((\bm\beta_{1,0},\bm\beta_{2,0})\in\Gamma_{\alpha}^{\mathrm{joint}}(\bm Y^{\obs})\right)
\ge \alpha-\delta_L^{\mathrm{partial}}.
\]
\end{proof}


\begin{theorem}[Marginal confidence set for $\bm{\beta}_1$]\label{thm:confidence set partial beta1}
Fix any $\alpha \in (0,1)$. Consider the partially permuted design~\eqref{eq:model_partial} 
defined in~\eqref{eq:confset-partial-beta1}. It holds that, for any $L \ge 2$,
\[
    \bP_{(\mathcal{U}^L, \bm{U})}\left( \bm{\beta}_1 \in \Gamma_{\alpha}^{\mathrm{partial}}(\bm{Y}^{\obs}) \right)
    \ge
    \alpha - \delta_L^{\mathrm{partial}},
\]
where $\delta_L^{\mathrm{partial}}$ is defined in~\eqref{eq:recovery-finite-L-partial-gammaL}.
\end{theorem}

\begin{proof}
By definition, we have
\begin{equation}\label{eq:confset-partial-beta1}
\Gamma_{\alpha}^{\mathrm{partial}}(\bm{Y}^{\obs})
:=\bigcup_{\Pi \in \mathcal{C}_{\mathrm{partial}}^{(L)}} \left\{
\bm\beta\in \R^{p_1}:\ 
T(\bm{Y}^{\obs}, (\Pi,\bm\beta) \mid \bm Z) \le F_{p_1, n-p}^{-1}(\alpha)
\right\}.
\end{equation}

Hence, we have
\[
\bP\left( \bm{\beta}_1 \in \Gamma_{\alpha}^{\mathrm{partial}}(\bm{Y}^{\obs}) \right)
\ge
\bP\left(
T\left(\bm{Y}^{\obs}, (\Pi_0, \bm{\beta}_1) \mid \bm{Z}\right) \le F_{p_1, n-p}^{-1}(\alpha)
\right)
-
\bP\left( \Pi_0 \notin \mathcal{C}_{\mathrm{partial}}^{(L)} \right).
\]
Under the true parameters $(\Pi_0, \bm{\beta}_1)$, the statistic
$T(\bm{Y}^{\obs}, (\Pi_0, \bm{\beta}_1)\mid \bm{Z})$ has an $F_{p_1, n-p}$ distribution, so the first probability equals $\alpha$.
Theorem~\ref{thm:recovery partial} bounds the second probability by $\delta_L^{\mathrm{partial}}$.
\end{proof}

\subsection{Theory on permuted ridge regression}\label{subsec:ridge_theory}

Recall the augmented design matrix $\bm X_\lambda:=\begin{pmatrix}\bm X\\ \sqrt{\lambda}I_p\end{pmatrix}\in\R^{(n+p)\times p}$ and, for any $\Pi\in\mP_{n,k}$, the embedded permutation $\tilde\Pi:=\mathrm{diag}(\Pi,I_p)$ and the corresponding augmented design
\[
\bm W_\lambda(\Pi):=\tilde\Pi\bm X_\lambda=\begin{pmatrix}\Pi\bm X\\ \sqrt{\lambda}I_p\end{pmatrix}.
\]
Let $P_{\Pi,\lambda}:=I_{n+p}-M_{\bm W_\lambda(\Pi)}$ be the residual projector.  Note that $\mathrm{rank}(P_{\Pi,\lambda})=n$.
For the augmented Gaussian perturbation $\tilde{\bm U}^*:=(\bm U^*,\bm 0_p)\in\R^{n+p}$ with $\bm U^*\sim\Normal(0,I_n)$,
we have $\tilde{\bm U}^*\sim\Normal(0,D)$ where $D:=\mathrm{diag}(I_n,0_p)$.

\begin{lemma}\label{lem:ridge_nondegenerate}
For any $\lambda>0$, define
\[
c_\lambda := \frac{\lambda}{\lambda+\|\bm X\|_{\rm op}^2}\in(0,1).
\]
Fix $\Pi\in\mP_{n,k}$ and define $Z_\Pi:=P_{\Pi,\lambda}\tilde{\bm U}^*$.
Let $\Sigma_\Pi$ denote the covariance of $Z_\Pi$ restricted to $\mathrm{range}(P_{\Pi,\lambda})$.
Then, it holds that
\[
c_\lambda I_n \preceq \Sigma_\Pi \preceq I_n.
\]
\end{lemma}

\begin{proof}
Set $P:=P_{\Pi,\lambda}$ and write $v=(a,b)\in\R^{n}\times\R^{p}$.
For any $v\in\mathrm{range}(P)$ we have $Pv=v$ and hence it holds that
\[
\mathrm{Var}(v^\top Z_\Pi)=v^\top P D P v = v^\top D v=\|a\|_2^2.
\]
It therefore suffices to lower and upper bound $\|a\|_2^2$ over unit vectors $v\in\mathrm{range}(P)$.

Let $v\in\mathrm{range}(P)$ with $\|v\|_2=1$.
Since $\mathrm{range}(P)=\mathrm{col}(\bm W_\lambda(\Pi))^\perp$, we have $\bm W_\lambda(\Pi)^\top v=0$, i.e.,
$\bm X^\top\Pi^\top a+\sqrt{\lambda}b=0$ and thus $b=-(1/\sqrt{\lambda})\bm X^\top\Pi^\top a$.
Consequently, we have
\[
1=\|v\|_2^2=\|a\|_2^2+\|b\|_2^2
=\|a\|_2^2+\frac{1}{\lambda}\|\bm X^\top\Pi^\top a\|_2^2
\le \|a\|_2^2\left(1+\frac{\|\bm X\|_{\rm op}^2}{\lambda}\right),
\]
which yields $\|a\|_2^2\ge \lambda/(\lambda+\|\bm X\|_{\rm op}^2)=c_\lambda$.
The upper bound $\|a\|_2^2\le 1$ is immediate from $\|v\|_2=1$.
Because of $\mathrm{Var}(v^\top Z_\Pi)=\|a\|_2^2$ for all unit $v\in\mathrm{range}(P)$, the claim follows.
\end{proof}

\begin{lemma}\label{lem:max_sim_ridge}
Fix $\Pi_0\in\mP_{n,k}$ and let $\tilde\Pi_0=\mathrm{diag}(\Pi_0,I_p)$.
For $\Pi\in\mP_{n,k}$ define $P_{\Pi,\lambda}:=I_{n+p}-M_{\bm W_\lambda(\Pi)}$ and consider $
\phi\left(P_{\Pi,\lambda}\tilde{\bm U}^*,P_{\Pi,\lambda}\tilde\Pi_0\bm X_\lambda\bm\beta_0\right)$.

Then for any $\gamma_1\in(0,1)$, it holds that
\[
\bP\left(
\max_{\Pi\in\mP_{n,k}}
\phi\left(P_{\Pi,\lambda}\tilde{\bm U}^*,P_{\Pi,\lambda}\tilde\Pi_0\bm X_\lambda\bm\beta_0\right)
\ge \gamma_1^2
\right)
\le
2|\mP_{n,k}|c_\lambda^{-n/2}\left(\arccos(\gamma_1)\right)^{n-1}.
\]
\end{lemma}

\begin{proof}
Fix $\Pi\in\mP_{n,k}$ and abbreviate $P:=P_{\Pi,\lambda}$.
If $P\tilde\Pi_0\bm X_\lambda \bm\beta_0=0$, then the corresponding $\phi(\cdot,\cdot)$ is zero and can be ignored.
Otherwise, set $w:=P\tilde\Pi_0\bm X_\lambda\bm\beta_0\neq 0$.

Let $Q\in\R^{(n+p)\times n}$ be any matrix with orthonormal columns spanning $\mathrm{range}(P)$.
Then, we have $Q^\top P Q=I_n$, $Q^\top w\neq 0$, and $\bar Z:=Q^\top P\tilde{\bm U}^*\sim\Normal(0,\Sigma)$ with
$\Sigma:=Q^\top D Q$.
Moreover, since $Q$ is an isometry on $\mathrm{range}(P)$, we have
\[
\phi(P\tilde{\bm U}^*,w)=\phi(\bar Z,\bar w),\qquad \bar w:=Q^\top w\in\R^{n}.
\]
By Lemma~\ref{lem:ridge_nondegenerate}, $\Sigma$ satisfies $c_\lambda I_n\preceq \Sigma\preceq I_n$ and hence
$\Sigma^{-1}$ has eigenvalues in $[1,1/c_\lambda]$.

Define $V:=\bar Z/\|\bar Z\|_2\in\mathbb{S}^{n-1}$ and $\bar e:=\bar w/\|\bar w\|_2$.
Then, it holds that $\phi(\bar Z,\bar w)=(V^\top \bar e)^2$, so we have
\[
\bP\left(\phi(P\tilde{\bm U}^*,w)\ge \gamma_1^2\right)=\bP\left(|V^\top \bar e|\ge \gamma_1\right).
\]
For $\bar Z\sim\Normal(0,\Sigma)$, the direction $V$ admits the density
\[
f(v)=\frac{(v^\top\Sigma^{-1}v)^{-n/2}}{\int_{\mathbb{S}^{n-1}}(u^\top\Sigma^{-1}u)^{-n/2}d\mu(u)}
\qquad (v\in\mathbb{S}^{n-1})
\]
with respect to the uniform surface measure $\mu$ on $\mathbb{S}^{n-1}$.
Since $(v^\top\Sigma^{-1}v)^{-n/2}\in[c_\lambda^{n/2},1]$ for all $v$, the normalizing constant is at least $c_\lambda^{n/2}$,
and hence $f(v)\le c_\lambda^{-n/2}$ for all $v$. Therefore, for any measurable $A\subseteq\mathbb{S}^{n-1}$, we have
$\bP(V\in A)\le c_\lambda^{-n/2}\mu(A)$.

Taking $A=\{v:|v^\top \bar e|\ge \gamma_1\}$ and using the bound
$\mu(A)\le 2(\arccos(\gamma_1))^{n-1}$ yields
\[
\bP\left(\phi(P\tilde{\bm U}^*,w)\ge \gamma_1^2\right)
\le 2c_\lambda^{-n/2}\left(\arccos(\gamma_1)\right)^{n-1}.
\]
A union bound over $\Pi\in\mP_{n,k}$ completes the proof.
\end{proof}

We next state the recovery guarantee for the ridge-augmented candidate set.
Define 
\begin{equation}\label{eq:c-min-ridge}
C_{\min}^{\mathrm{ridge}}(\Pi_0)
:=
\min_{\substack{\Pi \in \mP_{n, k},\\ \Pi \neq \Pi_0}}
\frac{\left\|(I_{n+p}-M_{\bm W_\lambda(\Pi)})\bm W_\lambda(\Pi_0)\bm\beta_0\right\|_2^2}
{n \max (d(\Pi_0, I_n) - d(\Pi, I_n), 1)}.
\end{equation}

\begin{theorem}\label{thm:candidate_set_ridge}
Assume $C_{\min}^{\mathrm{ridge}}(\Pi_0)>0$ holds and fix $\underline c\in(0,1)$.
Suppose $\lambda\ge \lambda_{\min}:=\|\bm X\|_{\rm op}^2\underline c/(1-\underline c)$ holds, so that $c_\lambda\ge \underline c$.
Then the conclusion of Theorem~\ref{thm:recovery} continues to hold for the confidence set $\mC_\lambda^{(L)}$,
with $C_{\min}(\Pi_0)$ replaced by $C_{\min}^{\mathrm{ridge}}(\Pi_0)$ and with the last term of $\Delta(\gamma)$
(see \eqref{eq:Delta-gamma}) similarly replaced by
\[
2^{(n+1)/2}|\mP_{n,k}|\underline c^{-n/2}\left(\gamma\log(e/\gamma)\right)^{(n-1)/2}.
\]
\end{theorem}

\begin{proof}
Fix $\gamma\in(0,1/4]$.
Recall the augmented perturbation $\tilde{\bm U}^*:=(\bm U^*,\bm 0_p)\sim\mathcal N(0,D)$ with
$D=\mathrm{diag}(I_n,0_p)\preceq I_{n+p}$, and the augmented designs
$\bm W_\lambda(\Pi)=\tilde\Pi\bm X_\lambda$ with residual projectors
$P_{\Pi,\lambda}:=I_{n+p}-M_{\bm W_\lambda(\Pi)}$ (so $\rank(P_{\Pi,\lambda})=n$).

We follow the proof of Lemma~\ref{lem:inconsistency} with $\bm W(\Pi)$ replaced by $\bm W_\lambda(\Pi)$,
$\bm U^*$ replaced by $\tilde{\bm U}^*$, and with the separation constant
$C_{\min}(\Pi_0)$ replaced by $C_{\min}^{\rm ridge}(\Pi_0)$.
Indeed, for any orthogonal projector $P$ of rank $n$ and any $x\ge 1$, write
$\tilde{\bm U}^*=D^{1/2}G$ with $G\sim\mathcal N(0,I_{n+p})$. Then, we have
\[
\tilde{\bm U}^{*\top}P\tilde{\bm U}^*
=G^\top\left(D^{1/2}PD^{1/2}\right)G,
\qquad
0\preceq D^{1/2}PD^{1/2}\preceq P,
\]
so all eigenvalues of $D^{1/2}PD^{1/2}$ lie in $[0,1]$. Hence
$\tilde{\bm U}^{*\top}P\tilde{\bm U}^*$ is stochastically dominated by $\chi_n^2$, and therefore we have
\[
\bP\left(\tilde{\bm U}^{*\top}P\tilde{\bm U}^*\ge nx\right)
\le \bP(\chi_n^2\ge nx)
\le \exp\left(-\frac n2\Psi(x)\right),
\]
which is exactly the tail bound used in Lemma~\ref{lem:inconsistency}.
Likewise, for any fixed $v$, $v^\top \tilde{\bm U}^*$ is Gaussian with variance $v^\top D v\le \|v\|_2^2$,
so the corresponding Gaussian tail bounds are also unchanged.

The only place in Lemma~\ref{lem:inconsistency} where spherical symmetry of projected perturbations is used
is in controlling the following event
\[
\max_{\Pi\in\mP_{n,k}}
\phi\left(P_{\Pi,\lambda}\tilde{\bm U}^*,P_{\Pi,\lambda}\tilde\Pi_0\bm X_\lambda\bm\beta_0\right).
\]
Here Lemma~\ref{lem:max-sim} is replaced by Lemma~\ref{lem:max_sim_ridge}, which yields the modified last
term $2^{(n+1)/2}|\mP_{n,k}|\underline c^{-n/2}(\gamma\log(e/\gamma))^{(n-1)/2}$.

Combining these bounds with the union bound over $L$ i.i.d. draws
$\tilde{\bm U}_1^*,\dots,\tilde{\bm U}_L^*$ yields the desired claim.
\end{proof}




\section{Additional Figures}\label{Appendix:figures}
In addition to Section \ref{sec:numerical}, we report the Monte Carlo averages of the coverage probability of the candidate set (Figure \ref{fig:app-mc-coverage}) and the volume of the coefficient confidence sets (Figure \ref{fig:app-coef-volume}).

\begin{figure}[t]
\centering
\includegraphics[width=0.95\linewidth]{figures/mc_coverage_all.png}
\caption{Monte Carlo averages of the coverage probability of the candidate set.}
\label{fig:app-mc-coverage}
\end{figure}


\begin{figure}[t]
\centering
\includegraphics[width=0.95\linewidth]{figures/coef_volume_all.png}
\caption{Monte Carlo averages of the volume of the confidence sets for the coefficient vector.}
\label{fig:app-coef-volume}
\end{figure}

\bibliographystyle{alpha}
\bibliography{mybib}



\end{document}
